\documentclass[11pt,reqno]{amsart}

%% Packages, typography and theorem environments.
%% The mathematical notation lives in notation.tex, the numerals in results.tex.

%% ------------------------------------------------------------------ packages
\usepackage[T1]{fontenc}
\usepackage[margin=1.15in]{geometry}
\usepackage{amsmath,amssymb,amsthm,mathtools}
\usepackage{booktabs,array}
\usepackage{microtype}
\usepackage{xcolor}
\definecolor{refblue}{RGB}{20,60,140}
\definecolor{todored}{RGB}{160,30,30}
\usepackage[colorlinks=true,linkcolor=refblue,citecolor=refblue,
            urlcolor=refblue,linktocpage=true,
            bookmarksnumbered=true,bookmarksdepth=2,
            pdftitle={Improved exponents for 3SUM and APSP from triangles in
              sparse lopsided graphs},
            pdfauthor={Jihoon Hyun},
            pdfkeywords={3SUM, APSP, Exact Triangle, thin matrix products,
              formalized mathematics}]{hyperref}

%% ---------------------------------------------------------------- typography
\allowdisplaybreaks
%% Lean names are set in a typewriter face and do not hyphenate.
\emergencystretch=3em

%% The parts of a statement are labeled (i), (ii), ...
\renewcommand{\theenumi}{\roman{enumi}}
\renewcommand{\labelenumi}{\textup{(\theenumi)}}

%% -------------------------------------------------------- Lean names
%% \lean{name} typesets the name of a Lean declaration in a typewriter face.
%% It is used every time a result is machine-checked.  The name is set by the
%% machinery of the url package (loaded by hyperref), so underscores, primes
%% and dots are safe and a long name may break after an underscore or a dot.
%% Do not use it in section titles or captions.  \leanraw is the unbreakable
%% form, \texttt{\detokenize{...}}.
\DeclareUrlCommand\leanurl{\urlstyle{tt}%
  \def\UrlBreaks{\do\.\do\_\do\/}\def\UrlBigBreaks{}%
  \def\UrlSpecials{\do\'{\mbox{\textquotesingle}}}}
\newcommand{\lean}[1]{\leanurl{#1}}
\newcommand{\leanraw}[1]{\texttt{\detokenize{#1}}}

%% -------------------------------------------------------- status markers
%% A numbered statement with no marker in its heading is machine-checked in
%% Lean, and its text names the declarations.  The one marker below goes into
%% the heading of the other statements, all of which are in Section 9:
%%   \nf    proved here on paper, not machine-checked --- never "unproved".
%% See the Formalization paragraph of the introduction.
\newcommand{\nf}{\textup{not formalized}}

%% A visible note for the parts of the draft that are still open.
\newcommand{\TODO}[1]{\textcolor{todored}{[\textsc{todo}: #1]}}

%% ------------------------------------------------------ theorem environments
\theoremstyle{plain}
\newtheorem{theorem}{Theorem}[section]
\numberwithin{equation}{section}
\numberwithin{table}{section}
\newtheorem{proposition}[theorem]{Proposition}
\newtheorem{lemma}[theorem]{Lemma}
\newtheorem{corollary}[theorem]{Corollary}

\theoremstyle{definition}
\newtheorem{definition}[theorem]{Definition}
\newtheorem{example}[theorem]{Example}

\theoremstyle{plain}
\newtheorem{question}[theorem]{Question}

\theoremstyle{remark}
\newtheorem{remark}[theorem]{Remark}

%% Mathematical notation.  Loaded after preamble.tex.

\newcommand{\R}{\mathbb{R}}
\newcommand{\N}{\mathbb{N}}
\newcommand{\Z}{\mathbb{Z}}
\newcommand{\Q}{\mathbb{Q}}

%% ------------------------------------------------------------ the recursion
\newcommand{\Hent}{H}                         % entropy, natural logarithms
\newcommand{\gtwo}{g_{2}}                     % exact rate of the leaf counts
\newcommand{\fone}{f_{1}}                     % rate of |W| alpha_d
\newcommand{\gX}{\gamma_{X}}                  % exact exponent of the boxes
\newcommand{\gpap}{\gamma_{\mathrm{pap}}}     % the paper's exponent
\newcommand{\gQ}{\gamma_{Q}}                  % exponent allowed by the queries
\newcommand{\Gam}{\Gamma}                     % common value at the crossing
\newcommand{\Rc}{R_{c}}                       % thinness bound, paper's encodings
\newcommand{\RcP}{R'_{c}}                     % thinness bound, pruned encodings
\newcommand{\Afull}{A_{\mathrm{full}}}        % thinness constant, all leaves
\newcommand{\Apr}{A_{\mathrm{pr}}}            % thinness constant, pruned
\newcommand{\Pz}{P_{0}}                       % the shared term of the identity
\newcommand{\binomial}[2]{\binom{#1}{#2}}
\DeclareMathOperator{\Leaves}{Leaves}
\DeclareMathOperator{\sav}{s}                 % the saving of a tuple

%% -------------------------------------------------------------- problems
\newcommand{\LopTri}{\textup{Lop-AE-SparseTri}}
\newcommand{\minplus}{(\min,+)}
\newcommand{\schon}[3]{\langle #1,#2,#3\rangle}
\newcommand{\polylog}{\operatorname{polylog}}

%% The numerals of the paper, one macro each.  They may still change while the
%% formalization is being completed; change them here and nowhere else.

%% ---------------------------------------------------------- the original paper
\newcommand{\ETpaper}{0.00175}            % saving for Exact Triangle (Thm 19)
\newcommand{\ETexpPaper}{2.9983}          % as printed: 3 - eps_T, eps_T = 0.0017
\newcommand{\ThreeSumPaper}{1.999125}     % 2 - \ETpaper/2
\newcommand{\ThreeSumPaperStated}{1.9992} % as printed in Theorem 22
\newcommand{\APSPpaper}{2.99942}          % 3 - \ETpaper/3, as printed

%% --------------------------------- T2: the paper's exponents, better parameters
\newcommand{\dTwo}{0.00184}
\newcommand{\ThreeSumTwo}{1.99908}
\newcommand{\ETTwo}{2.99816}
\newcommand{\APSPTwo}{2.99939}

%% ------------------------------------------------------------ T3: + exact count
\newcommand{\dThree}{0.00196}
\newcommand{\ThreeSumThree}{1.99902}
\newcommand{\ETThree}{2.99804}
\newcommand{\APSPThree}{2.99935}

%% ------------------------------------------------------ T4: + true middle size
\newcommand{\dFour}{0.00205}
\newcommand{\ThreeSumFour}{1.99898}
\newcommand{\ETFour}{2.99795}
\newcommand{\APSPFour}{2.99932}

%% --------------------------------------------- T1: + pruned encodings
\newcommand{\dOne}{0.002095}
\newcommand{\ThreeSumOne}{1.99896}
\newcommand{\ETOne}{2.99791}
\newcommand{\APSPOne}{2.99931}
\newcommand{\APSPFn}{2.99795}             % (min,+), APSP through all-edges ET, no hypothesis
\newcommand{\APSPFnOne}{2.99791}          % the same with pruned encodings (T1)

%% ----------------------------------------------------------------- the optimum
\newcommand{\SstarLo}{0.002095}           % lower bound for sup_c S(c), pruned
\newcommand{\SstarHi}{0.002096}           % upper bound for sup_c S(c), pruned
\newcommand{\cStarLo}{21}                 % S(c) >= \SstarLo only for c in
\newcommand{\cStarHi}{23}                 %   (\cStarLo, \cStarHi)
\newcommand{\cUncond}{108/5}              % ratio c of Theorems T2, T3
\newcommand{\cUncondOpt}{107/5}           % ratio c of Theorem T4
\newcommand{\cCond}{22}                   % ratio c of Theorem T1 (pruned encodings)
\newcommand{\SfullLo}{0.002059}           % lower bound for sup_c S(c), paper's encodings
\newcommand{\SfullHi}{0.002061}           % upper bound for sup_c S(c), paper's encodings
\newcommand{\cFullLo}{20.8}               % S_full(c) >= \SfullLo only for c in
\newcommand{\cFullHi}{22}                 %   (\cFullLo, \cFullHi)

%% approximate optimum with pruned encodings (not certified)
\newcommand{\SprApprox}{0.0020953}
\newcommand{\cPrApprox}{21.92}
\newcommand{\GamPrApprox}{0.0759}
\newcommand{\sigmaPrApprox}{0.519}
\newcommand{\epsPrApprox}{0.0552}
\newcommand{\thetaPrApprox}{0.116}
%% approximate optimum with the paper's encodings (not certified)
\newcommand{\SfullApprox}{0.0020599}
\newcommand{\cFullApprox}{21.38}
\newcommand{\cFitsMax}{21.86}             % the tile fits for c below this
\newcommand{\cFullMin}{21.49}             % A_full(c) >= 18 ln 4 for c above this

%% parameters of Theorem T1 at c = \cCond
\newcommand{\GamCond}{0.0762}
\newcommand{\sigmaCond}{0.519}
\newcommand{\epsCond}{0.05498}
\newcommand{\thetaCond}{0.116}

%% ------------------------------------------------------------ the Lean sources
\newcommand{\LeanFiles}{179}              % files of Lean in ImprovedExponents/
\newcommand{\LeanLines}{46\,416}          % lines of Lean in ImprovedExponents/
\newcommand{\LeanCellLines}{26\,822}      % of which generated (GlobalCells*, GlobalFullCells*, Shapes/Cells_*)
\newcommand{\LeanHandLines}{19\,600}      % hand-written, about
\newcommand{\LeanExactCountLines}{763}    % lines of ImprovedExponents/ExactCount/
\newcommand{\LeanPipelineLines}{2337}     % lines of ImprovedExponents/Pipeline/
\newcommand{\LeanOptimumLines}{29\,754}   % lines of ImprovedExponents/Optimum/
\newcommand{\LeanGlobalCellLines}{4703}   % of which GlobalCells*, GlobalFullCells*
\newcommand{\LeanShapesHandLines}{1038}   % hand-written lines of Optimum/Shapes/
\newcommand{\LeanShapesCellLines}{22\,119} % generated lines of Optimum/Shapes/Cells_*
\newcommand{\LeanShapesCells}{2253}       % certified cells of the nine shapes
\newcommand{\LeanShapesCellFiles}{53}     % files holding them
\newcommand{\LeanPaperAccountingLines}{269} % lines of Optimum/PaperAccounting.lean
\newcommand{\LeanPrunedProgramLines}{3951} % lines of ImprovedExponents/PrunedProgram/
\newcommand{\LeanCostEightPLines}{718}    % lines of ImprovedExponents/Cost8P/
\newcommand{\LeanStatementsLines}{396}    % lines of Statements.lean
\newcommand{\LeanAllEdgesLines}{3558}    % lines of ImprovedExponents/AllEdges/
\newcommand{\LeanCostEightXLines}{727}    % lines of ImprovedExponents/Cost8X/
\newcommand{\LeanPrunedEncodingLines}{1330} % lines of ImprovedExponents/PrunedEncoding/
\newcommand{\LeanHostMidLines}{1008}      % lines of ImprovedExponents/HostMid/
\newcommand{\LeanParamRoutinesLines}{645} % lines of ImprovedExponents/ParamRoutines/
\newcommand{\LeanTotalLines}{806}         % lines of ImprovedExponents/Total/
\newcommand{\LeanMinPlusLines}{400}       % lines of ImprovedExponents/MinPlus/
\newcommand{\LeanVersion}{v4.33.1}        % Lean and Mathlib of this project (lean-toolchain), as upstream
\newcommand{\ComparatorDate}{2026-10-07}  % the day Comparator accepted comparator.json (scripts/comparator-run.log)
\newcommand{\UpstreamRev}{e1a4e65}        % revision of the upstream repository
\newcommand{\EncConst}{778}               % the constant of tEncP_le


\title[Improved exponents for 3SUM and APSP]{Improved exponents for 3SUM and
APSP\\ from triangles in sparse lopsided graphs}

\author{Jihoon Hyun}
%% \address{...}
\email{qawbecrdtey@kaist.ac.kr}

\subjclass[2020]{Primary 68Q25; Secondary 68W40, 68V20}

\keywords{3SUM, all-pairs shortest paths, Exact Triangle, thin matrix
products, Sch\"onhage's identity, formalized mathematics, Lean}

\date{\today}

\begin{document}

\begin{abstract}
Alman and Vassilevska Williams gave deterministic algorithms for 3SUM in time
\(O(n^{\ThreeSumPaperStated})\) and for the \((\min,+)\)-product and APSP in
time \(O(n^{\APSPpaper})\), through an algorithm for Exact Triangle in time
\(n^{3-\ETpaper}\).  We analyze their algorithm again and change it in a few
places.  An exact count of the leaves of the recursion, a host that charges the
solver at the true size of the middle part, and encodings restricted to the
leaves that are read raise the saving for Exact Triangle from \(\ETpaper\) to
every number below \(\SstarLo\), so that 3SUM is solved in time
\(O(n^{\ThreeSumOne})\); we show that the supremum of the savings of the method
is at most \(\SstarHi\).  Following a footnote of the paper, we also solve the
all-edges version of Exact Triangle within the same bound and reduce the
\((\min,+)\)-product to it on the same number of vertices, so that APSP keeps
the whole saving: the \((\min,+)\)-product and APSP are solved in time
\(O(n^{\APSPFnOne})\), against \(O(n^{\APSPOne})\) through the standard
reductions.  All statements are about programs of a word RAM, hold without any
hypothesis, and are machine-checked in Lean~4 on top of the formalization of
the original paper; the pruned encoder is written as a program and its running
time is proved on the machine.
\end{abstract}

\maketitle

\setcounter{tocdepth}{1}
\tableofcontents

\section{Introduction}
\label{sec:intro}

Alman and Vassilevska Williams \cite{AV} gave the first truly subquadratic
algorithm for 3SUM and the first truly subcubic algorithm for APSP.  Both
follow from a deterministic algorithm for Exact Triangle in time
\(n^{3-\ETpaper}\) up to a logarithmic factor, which rests on a fast algorithm
for \emph{thin} matrix products: a recursion with Sch\"onhage's identity for
\(\schon313\oplus\schon141\), pruned to the entries that are wanted.  The
savings are small, and the paper does not optimize them.  The result has since
been formalized in Lean~4 as a set of statements about programs of a word RAM
\cite{upstream}.

This paper asks how far the same method goes with the same base identity.  It
has three ingredients, a determination of the optimal parameters, and the
all-edges route to APSP that footnote~10 of the paper mentions.

\subsection*{The ingredients}

\emph{An exact count of the leaves} (Section~\ref{sec:exact}).  The cost of the
recursion is governed by the numbers \(\beta_d\) of leaves of each order in a
tile.  The paper bounds their tail by a geometric series.  We bound it by its
exponential rate, a function \(\gtwo\) of the ratio \(c=L/m\) of the levels and
of the relative order (Theorem~\ref{thm:tail}); the proof is a comparison of
integers.  The paper's exponent is the tangent of the exact one at zero.  The
algorithm does not change: the same program gets a better bound on its running
time (Theorem~\ref{thm:offline}).

\emph{The true size of the middle part} (Section~\ref{sec:middle}).  The
reduction from Exact Triangle (Theorem~17 of the paper) produces instances
whose middle part has about \(D'/g'\) vertices, where \(D'\) is the dimension
parameter and \(g'\) the number of pieces per hash class, and calls the solver
as if it had \(D'\).  Passing the true size changes one assignment of the
program (Theorem~\ref{thm:mid}).  The solver then runs at a smaller inner
dimension \(D\) and, relative to it, with fewer wanted entries:
\(n^{2}/D^{\sigma}\) with \(\sigma>1/2\) instead of \(n^{2}/\sqrt D\).
Equivalently, the weights are hashed modulo a larger prime
(Remark~\ref{rem:prime}).

\emph{Pruned encodings} (Section~\ref{sec:pruned}).  The recursion reads the
encodings of the input only at the leaves with at most \(m\) symbols \(\Pz\).
They can be computed by a pruned form of Yates' algorithm in \(O(LM)\)
operations per band instead of \(O(10^{L})\) (Theorem~\ref{thm:pruned}), which
relaxes the condition under which a product counts as thin.  We write the
pruned encoder as a program and prove its running time on the machine
(Theorem~\ref{thm:encoder}); the rest of the solver is upstream's, re-verified
under the weaker contract that the encodings are right only at the leaves that
are read.

\emph{APSP through all-edges Exact Triangle} (Section~\ref{sec:alledges}).
The known reductions give the \(\minplus\)-product and APSP a third of the
saving of Exact Triangle.  Footnote~10 of the paper observes that the host of
its reduction also solves the all-edges version of Exact Triangle, to which the
\(\minplus\)-product reduces on the same number of vertices.  We write the two
programs (the all-edges host and the reduction, Theorems~\ref{thm:mid-ae}
and~\ref{thm:pairs}), and APSP keeps the whole saving.

\subsection*{The optimum}

With these ingredients the saving for Exact Triangle depends on four
parameters.  For a fixed ratio \(c\) its supremum has a closed form: it is
\(R(\Gam)\Gam/2\), where \(\Gam=\Gam(c)\) is the value at which two explicit
curves cross and \(R\) is the bound of the thinness condition
(Theorem~\ref{thm:sup}).  With pruned encodings the supremum over all ratios
lies in the interval \((\SstarLo,\SstarHi]\), and only ratios in
\((\cStarLo,\cStarHi)\) come within \(\SstarLo\)
(Theorem~\ref{thm:global}); with the paper's encodings it lies in
\((\SfullLo,\SfullHi]\), approached only for ratios in \((\cFullLo,\cFullHi)\)
(Theorem~\ref{thm:global-full}).  The theorems reach every saving below the
two suprema: the upstream program tests \(D^{18}\le N\) to decide whether its
recursion applies, which excludes the optimal ratios, and we replace the test
by one with a rational exponent (Section~\ref{sec:regime}).
So the method, with this identity and these ingredients, ends at a saving of
about \(\SprApprox\) for Exact Triangle, that is, at the exponent
\(2-\SprApprox/2\) for 3SUM.  Section~\ref{sec:limits} explains why no identity
of the same kind does better.

\subsection*{Results}

Table~\ref{tab:results} lists the exponents.  The column ``paper'' is
\cite{AV}; each of the next columns adds one ingredient to the previous one.

\begin{table}[ht]
\centering\small
\setlength{\tabcolsep}{5pt}
\begin{tabular}{lccccc}
\toprule
 & paper
 & Thm.~\ref{thm:T2}
 & Thm.~\ref{thm:T3}
 & Thm.~\ref{thm:T4}
 & Thm.~\ref{thm:T1}\\
\midrule
added & & parameters & exact count & middle size & pruned enc.\\
saving \(\delta\)
 & \(\ETpaper\) & \(\dTwo\) & \(\dThree\) & \(\SfullLo\) & \(\dOne\)\\
Exact Triangle
 & \(\ETexpPaper\) & \(\ETTwo\) & \(\ETThree\) & \(\ETFour\) & \(\ETOne\)\\
3SUM
 & \(\ThreeSumPaperStated\) & \(\ThreeSumTwo\) & \(\ThreeSumThree\)
 & \(\ThreeSumFour\) & \(\ThreeSumOne\)\\
\(\minplus\)-product, APSP
 & \(\APSPpaper\) & \(\APSPTwo\) & \(\APSPThree\) & \(\APSPFour\) & \(\APSPOne\)\\
\quad all-edges route
 & & & & \(\APSPFn\) & \(\APSPFnOne\)\\
\midrule
status & \cite{upstream} & checked & checked & checked & checked\\
\bottomrule
\end{tabular}
\caption{Exponents \(a\) of the running times \(O(n^{a})\) of deterministic
word-RAM programs.  The exponents of the column ``paper'' are those printed in
\cite{AV} and proved in \cite{upstream}.  The all-edges route
(Theorem~\ref{thm:apsp-ae}) is carried out for the two solvers of the last two
columns.}
\label{tab:results}
\end{table}

All four columns of new results hold without any hypothesis.  The first of
them uses none of the ingredients: it is the algorithm and the analysis of the
paper at better parameters.  The first three use the existing encoder, which is
the paper's; the last uses the pruned encoder of Theorem~\ref{thm:encoder}.
Theorems~\ref{thm:T4} and~\ref{thm:T1} give every saving below the two suprema
of Theorems~\ref{thm:global-full} and~\ref{thm:global}; the decimal exponents
are their values at one ratio each.

\subsection*{Formalization}

The results are formalized in Lean~4 \cite{Lean4} with Mathlib \cite{Mathlib},
on top of the formalization \cite{upstream} of the original paper; the
development and this paper are in the repository
\url{https://github.com/qawbecrdtey/ImprovedExponents}.  Every time
a result is machine-checked, the text names the Lean declarations in a
typewriter face, as in \lean{sum_beta_le_exp}.  The heading of a numbered
statement shows its status.
\begin{itemize}
\item \emph{No marker}: the statement is machine-checked.
\item ``\nf'': the statement is proved here on paper and is not
  machine-checked.  The marker never means ``unproved''.  It occurs only in
  Section~\ref{sec:limits}, for the mixtures of shapes, the hypothetical
  shapes, the experiments and a remark on \(\omega\); every numbered statement
  outside Section~\ref{sec:limits} is machine-checked, and so is the
  comparison of the shapes in Section~\ref{sec:limits}.
\end{itemize}
Section~\ref{sec:formalization} says what has to be trusted and how the
development is checked.

\subsection*{Organization}

Section~\ref{sec:framework} recalls the framework of \cite{AV}.
Sections~\ref{sec:exact}, \ref{sec:pruned} and~\ref{sec:middle} have the three
ingredients, Section~\ref{sec:optimum} the optimal parameters and
Section~\ref{sec:consequences} the theorems for 3SUM, Exact Triangle, the
\(\minplus\)-product and APSP, together with the all-edges route to APSP.
Section~\ref{sec:formalization} describes the formalization and
Section~\ref{sec:limits} the limits of the method.

\section{The framework of Alman and Vassilevska Williams}
\label{sec:framework}

We recall the parts of \cite{AV} that are used.  Numbers of lemmas and theorems
``of the paper'' are those of \cite{AV}.  All logarithms are natural unless a
base is written, \(\Hent(x)=-x\ln x-(1-x)\ln(1-x)\) is the entropy function, and
\(\psi(x)=x\ln x\) with \(\psi(0)=0\).

\subsection{Problems and the machine}
\label{sec:machine}

\emph{Exact Triangle} asks whether a weighted tripartite graph with \(n\)
vertices per part has a triangle of weight zero.  \emph{3SUM} asks whether \(n\)
integers contain three with sum zero.  The \emph{\(\minplus\)-product} of two
\(n\times n\) integer matrices \(A,B\) is \(C_{ij}=\min_k(A_{ik}+B_{kj})\), and
\emph{APSP} asks for all distances in a directed graph with integer weights and
no negative cycle.

The \emph{thin product} is the following task.  Given \(X\in\Z^{N\times D}\),
\(Y\in\Z^{D\times N}\) and a set \(W\) of \emph{wanted} positions of the
\(N\times N\) matrix \(XY\), compute \((XY)[I,J]\) for all \((I,J)\in W\).  The
number \(D\) is the \emph{inner dimension}.  \(\LopTri(n,D)\)
(Definition~13 of the paper) is the special case that Exact Triangle reduces
to: a tripartite graph with outer parts of \(n\) vertices and a middle part of
at most \(D\) vertices, and a set \(W\) of \emph{query pairs} of outer
vertices; decide for each pair whether it has a common neighbor in the middle
part.

Running times are those of programs of a word RAM, in the sense of the
formalization \cite{upstream} of the paper.  For a problem \(Q\) and a rational
number \(r\), the statement \lean{Q.SolvedInTime}~\(r\) says: for every \(\kappa\)
there are a program, a slope \(b\) and a bound \(T(n)=O(n^{r})\) such that the
program is correct, within \(T(n)\) steps, on all instances of size \(n\) whose
numbers have absolute value at most \(n^{\kappa}\), at every word size
\(\ge b(\lfloor\log_2 n\rfloor+1)\).  The formalization proves Theorem~22 of
the paper in this form, with \(r=\ThreeSumPaperStated\) for 3SUM and
\(r=\APSPpaper\) for the \(\minplus\)-product and for APSP.

\subsection{The recursion}
\label{sec:recursion}

The recursion of Section~2 of the paper applies Sch\"onhage's identity
\cite{Sch81} for \(\schon313\oplus\schon141\).  The identity has ten terms.
Nine of them feed one outer output each, and the tenth, \(\Pz\), is shared.  It
has seven left variables: three outer ones and four inner ones.

The identity is applied on \(L\) levels, \(m\) of which are inner.  We write
\[
  D=4^{m},\qquad N_0=3^{L-m},\qquad K=\binomial Lm,\qquad
  M=K N_0^{2}=\binomial Lm 9^{L-m},\qquad c=\frac Lm\ge 10 .
\]
The matrix \(XY\) is cut into \emph{tiles} of side \(\sqrt K N_0\) (at most
\(4N^2/M\) of them), each with \(M\) output strings, and the rows of \(X\) into
\emph{bands} of \(\sqrt K N_0\) rows.

A \emph{leaf} is a string of \(L\) terms.  Its \emph{order} is \(m\) minus its
number of symbols \(\Pz\).  For a set \(U\) of output strings, \(\Leaves(U)\)
is the set of leaves at which the recursion multiplies for some string of
\(U\); these leaves have order between \(0\) and \(m\).  Two counts govern the
cost:
\[
  \alpha_d=\binomial md 9^{d},\qquad
  \beta_d=\binomial L{m-d}9^{L-m+d}\qquad(0\le d\le m).
\]
One output string is fed by \(\alpha_d\) leaves of order \(d\), and a tile has
\(\beta_d\) leaves of order \(d\); \(\beta_0=M\).  Lemma~11 of the paper is the
bound
\begin{equation}\label{eq:lemma11}
  |\Leaves(U)|\le\sum_{d=0}^{m}\min\bigl(|U|\alpha_d,\ \beta_d\bigr)
\end{equation}
for a set \(U\) of output strings of one tile, followed by an estimate of the
right-hand side.

\subsection{Boxes and queries}
\label{sec:boxes}

Section~4 of the paper (Theorem~30) turns the recursion into a data structure
with a switching order \(0\le t\le m\).  The leaves of order at least \(t\) are
covered by precomputed \emph{boxes}, at most
\begin{equation}\label{eq:boxes}
  (m+1)\sum_{d=t}^{m}\beta_d
\end{equation}
per tile (Lemma~29 of the paper).  A query for one entry sums over the leaves
of order below \(t\) and costs \(O\bigl(L\sum_{d\le t}\alpha_d\bigr)\).  The
encodings \(\Phi_\tau(a)\) of the input at all \(10^{L}\) leaves \(\tau\) are
computed for every band by Yates' algorithm \cite{Yat37}.  The preprocessing
costs a constant times the expression~(8) of the paper,
\begin{equation}\label{eq:cost8}
  L\,m\,\frac{\rho^{t}}{1-\rho}\,N^{2}+\frac{N\cdot 10^{L}}{\sqrt K N_0},
  \qquad \rho=\frac{9m}{L-m+1},
\end{equation}
because the paper bounds \(\sum_{d\ge t}\beta_d\le M\rho^{t}/(1-\rho)\); the
second summand is the cost of the encodings.

Fix real numbers \(c>10\) and \(0<\theta<0.9\) and let \(L=\lceil cm\rceil\),
\(t=\lceil\theta m\rceil\).  Then \(\rho<\rho_c:=9/(c-1)<1\).  Put
\begin{equation}\label{eq:gpap-q}
  \gpap(c,\theta)=\frac{\theta\ln(1/\rho_c)}{\ln 4},\qquad
  q(\theta)=\frac{\Hent(\theta)+\theta\ln 9}{\ln 4}.
\end{equation}
The first summand of \eqref{eq:cost8} is \(O(N^{2}\log^{2}D/D^{\gpap})\) and a
query costs \(O(D^{q(\theta)}\log D)\).  The second summand is at most
\(N^{2}/D^{\gamma}\) when the instance is \emph{thin}: \(D\le N^{\varepsilon}\)
with
\begin{equation}\label{eq:Rc}
  \varepsilon<\Rc(\gamma):=\frac{\ln 4}{\Afull(c)+\gamma\ln 4},\qquad
  \Afull(c)=c\ln 10-\frac c2\Hent\Bigl(\frac1c\Bigr)-(c-1)\ln 3 .
\end{equation}
This is equation~(11) of the paper; \(\Afull(c)\) is the exponential rate of
\(10^{L}/(\sqrt K N_0)\) per inner level.  Corollary~31 of the paper is the
resulting data structure, and Corollary~32 the offline form: for
\(|W|\le N^{2}/D^{\kappa}\) wanted entries the thin product takes time
\begin{equation}\label{eq:cor32}
  O\bigl(N^{2}\log^{2}D\,(D^{-\gamma}+D^{q(\theta)-\kappa})\bigr),
  \qquad\gamma=\gpap(c,\theta).
\end{equation}
Corollary~26 of the paper takes \(c=21\), \(\theta=1/9\), \(\kappa=1/2\) and
\(\varepsilon=1/18\), which gives \(\gamma=0.063\).

\subsection{From the thin product to Exact Triangle}
\label{sec:host}

Theorem~17 of the paper reduces Exact Triangle on \(n\) vertices per part with
weights of absolute value at most \(n^{\nu}\) to \(\LopTri\).  Let
\(16\le D'\le n\) and let \(1\le g'\le\sqrt{D'}\) be an integer.  The weights
are hashed modulo a prime \(p\in[\sqrt{D'}/2,\sqrt{D'})\) that is chosen to
minimize the number of false positives.  The reduction produces at most
\(4ng'\) instances of \(\LopTri(n,D')\) with at most \(n^{2}/\sqrt{D'}\) query
pairs each, and takes
\begin{equation}\label{eq:thm17}
  O\Bigl(\frac{\nu n^{3}\log n}{g'}+n^{\omega+o(1)}D'^{3/2}+n^{2}D'g'\Bigr)
\end{equation}
additional time.  The three terms are the scans for witnesses, the choice of the
prime and the construction of the instances.  The middle part of an instance is
\(C_k\times\Z_p\) for a piece \(C_k\) of one part with
\(|C_k|\le\lceil s/g'\rceil\), \(s=\lfloor\sqrt{D'}\rfloor\).  We write \(D'\)
and \(g'\) for the parameters of the reduction and keep \(D\) for the inner
dimension of the thin products; in the paper the two agree.

With \(D'=n^{1/18}\) and \(g'=D'^{\gamma/2}\), \(\gamma=0.063\), the calls and
the scans balance (Remark~20 of the paper), and Exact Triangle takes time
\(n^{3-\ETpaper}\) up to a logarithmic factor (Theorem~19 of the paper), since
\(0.063/36=\ETpaper\).

\subsection{From Exact Triangle to 3SUM and APSP}
\label{sec:known}

Theorem~21 of the paper recalls two known deterministic reductions.  3SUM on
\(n\) numbers reduces to \(n^{1/2+o(1)}\) instances of Exact Triangle on
\(n^{1/2+o(1)}\) vertices per part \cite{CH20,VW13}.  The \(\minplus\)-product
of \(n\times n\) matrices with entries up to \(U\) takes time
\(O(n^{2}T(n^{1/3})\log^{2}U)\) if Exact Triangle on \(s\) vertices per part
takes time \(T(s)\), and APSP follows with one more logarithm
\cite{VW10,VW18,VW13}.  So if Exact Triangle is solved in time \(n^{3-\delta}\),
then 3SUM is solved in time \(n^{a}\) for every \(a>2-\delta/2\), and the
\(\minplus\)-product and APSP in time \(n^{a}\) for every \(a>3-\delta/3\):
3SUM keeps half of the saving and APSP a third.  With \(\delta=\ETpaper\) this
is Theorem~22 of the paper: 3SUM in time \(n^{\ThreeSumPaper+o(1)}\), printed
as \(O(n^{\ThreeSumPaperStated})\), and APSP in time \(O(n^{\APSPpaper})\).

\section{The exact count of the leaves}
\label{sec:exact}

The paper bounds the tail \(\sum_{d\ge t}\beta_d\) of the leaf counts by a
geometric series.  This section bounds it by its true exponential rate.  The
program does not change; only its analysis does.

\subsection{The rate function}

\begin{definition}\label{def:g2}
For \(c>1\) and \(0\le\delta\le1\) let
\[
  \gtwo(c,\delta)=\psi(c-1)-\psi(c-1+\delta)-\psi(1-\delta)+\delta\ln 9
  =c\,\Hent\Bigl(\frac{1-\delta}{c}\Bigr)-c\,\Hent\Bigl(\frac1c\Bigr)+\delta\ln 9,
\]
and \(\gX(c,\theta)=-\gtwo(c,\theta)/\ln 4\).
\end{definition}

The two expressions agree because
\(c\,\Hent(x/c)=c\ln c-\psi(x)-\psi(c-x)\) for \(0\le x\le c\).  In Lean the
functions are \lean{g2} and \lean{gammaX}.

\begin{lemma}\label{lem:g2}
Let \(c\ge 10\).
\begin{enumerate}
\item \(\partial\gtwo/\partial\delta=\ln\bigl(9(1-\delta)/(c-1+\delta)\bigr)\le0\)
  for \(0<\delta<1\).  So \(\gtwo(c,\cdot)\) is nonincreasing on \([0,1]\), and
  strictly decreasing if \(c>10\).
\item \(\partial\gtwo/\partial c=\ln\bigl((c-1)/(c-1+\delta)\bigr)\le0\).  So
  \(\gtwo(\cdot,\delta)\) is nonincreasing on \([1,\infty)\).
\item \(\gtwo(c,0)=0\), and \(\gtwo(c,\delta)<0\) for \(c>10\) and
  \(0<\delta\le1\).
\item \(\gtwo(c,\delta)\le\delta\ln\rho_c\) with \(\rho_c=9/(c-1)\).  Hence
  \(\gpap(c,\theta)\le\gX(c,\theta)\) for \(0\le\theta\le1\), and for \(c>10\)
  the inequality is strict unless \(\theta=0\).
\end{enumerate}
\end{lemma}

\begin{proof}
(i) and (ii) follow from \(\psi'(x)=\ln x+1\).  The sign in (i) is that of
\(9(1-\delta)-(c-1+\delta)=10-c-10\delta\).  (iii) follows from (i).  For (iv),
the second derivative of \(\gtwo(c,\cdot)\) is
\(-1/(c-1+\delta)-1/(1-\delta)<0\), so the function is strictly concave and lies
below its tangent at \(0\), which is \(\delta\mapsto\delta\ln\rho_c\).
\end{proof}

The parts of the lemma are \lean{hasDerivAt_g2_right},
\lean{g2_antitoneOn_right}, \lean{g2_strictAntiOn_right},
\lean{hasDerivAt_g2_left}, \lean{g2_antitoneOn_left}, \lean{g2_neg} and
\lean{gammaOf_le_gammaX}; the strict inequalities of (iv), \(\gtwo(c,\delta)<\delta\ln\rho_c\)
for \(\delta>0\) and \(\gpap(c,\theta)<\gX(c,\theta)\) for \(\theta>0\) and
\(c>10\), are \lean{g2_lt_mul_log_rhoC} and \lean{gammaOf_lt_gammaX}.

\subsection{The tail}

\begin{theorem}[Exact tail]\label{thm:tail}
Let \(m\ge1\), \(L\ge10m\) and \(0\le t\le m\).  Then
\[
  \sum_{d=t}^{m}\beta_d\le(m+1)(L+1)\,M\,e^{m\,\gtwo(L/m,\,t/m)} .
\]
\end{theorem}

\begin{proof}
First,
\begin{equation}\label{eq:exp-g2}
  e^{m\,\gtwo(L/m,\,t/m)}
  =\frac{9^{t}\,m^{m}\,(L-m)^{L-m}}{(m-t)^{m-t}\,(L-m+t)^{L-m+t}},
\end{equation}
with \(0^{0}=1\).  Indeed \(m\,\psi(x/m)=\psi(x)-x\ln m\), so
\(m\,\gtwo(L/m,t/m)=\psi(L-m)+\psi(m)-\psi(L-m+t)-\psi(m-t)+t\ln 9\), the
multiples of \(\ln m\) cancelling.

So it suffices to prove, for \(t\le d\le m\), the inequality between integers
\begin{equation}\label{eq:beta-int}
  \beta_d\,(m-t)^{m-t}(L-m+t)^{L-m+t}\le(L+1)\,M\,9^{t}\,m^{m}(L-m)^{L-m},
\end{equation}
and to sum it over the at most \(m+1\) values of \(d\).  It follows from three
facts.
\begin{enumerate}
\item For \(j\le j_0\le L/10\),
  \(\binomial Lj9^{j_0-j}\le\binomial L{j_0}\): the ratio of consecutive
  binomial coefficients is \((L-j)/(j+1)\ge9\) for \(j<L/10\).
\item \(\binomial L{j_0}j_0^{j_0}(L-j_0)^{L-j_0}\le L^{L}\): the left-hand side
  is one term of the binomial expansion of \((j_0+(L-j_0))^{L}\).
\item \(L^{L}\le(L+1)\binomial Lm m^{m}(L-m)^{L-m}\): the expansion of
  \((m+(L-m))^{L}\) has \(L+1\) terms, and the term with index \(m\) is the
  largest.
\end{enumerate}
Apply (i) with \(j=m-d\) and \(j_0=m-t\), which is at most \(L/10\):
\(\beta_d=\binomial L{m-d}9^{d-t}\cdot9^{L-m+t}\le\binomial L{m-t}9^{L-m+t}\).
By (ii) with the same \(j_0\) and then (iii), the left-hand side of
\eqref{eq:beta-int} is at most
\(9^{L-m+t}L^{L}\le9^{L-m+t}(L+1)\binomial Lm m^{m}(L-m)^{L-m}\), which is the
right-hand side because \(M=\binomial Lm9^{L-m}\).
\end{proof}

In Lean: \lean{exp_mul_g2_div} is \eqref{eq:exp-g2},
\lean{choose_mul_nine_pow_le} is fact~(i), \lean{beta_mul_le} is
\eqref{eq:beta-int}, and \lean{sum_beta_le_exp} is the theorem.

\begin{corollary}\label{cor:tail}
Let \(c>10\), \(0\le\theta\le1\), \(m\ge1\), \(L=\lceil cm\rceil\),
\(t=\lceil\theta m\rceil\) and \(D=4^{m}\).  Then
\[
  \sum_{d=t}^{m}\beta_d\le(m+1)(L+1)\,M\,D^{-\gX(c,\theta)} .
\]
\end{corollary}

\begin{proof}
\(L/m\ge c\) and \(\theta\le t/m\le1\), so
\(\gtwo(L/m,t/m)\le\gtwo(c,\theta)\) by Lemma~\ref{lem:g2}(i),(ii), and
\(e^{m\,\gtwo(c,\theta)}=D^{-\gX(c,\theta)}\).
\end{proof}

In Lean: \lean{sum_beta_le_rpow}.

\subsection{The offline routine with the exact exponent}

By \eqref{eq:boxes} and Corollary~\ref{cor:tail}, the boxes of all tiles number
at most \(4(m+1)^{2}(L+1)N^{2}D^{-\gX(c,\theta)}\).  For \(\gamma<\gX(c,\theta)\)
the factor \((m+1)(L+1)\) is absorbed by \(D^{\gX-\gamma}\).  So the first
summand of \eqref{eq:cost8} can be replaced by \(O(N^{2}\log^{2}D/D^{\gamma})\)
for every \(\gamma<\gX(c,\theta)\); the expression with the exact tail in place
of the geometric bound is \lean{cost8X}.

\begin{theorem}[The offline routine]\label{thm:offline}
Let \(c=a/b>10\) and \(\theta=p/q\in(0,0.9)\) be rational, let
\(0<\gamma<\gX(c,\theta)\) and \(\varepsilon<\Rc(\gamma)\).  For every large
enough threshold \(m_0\), the offline routine of Section~4.4 of the paper at
the parameters \((c,\theta,m_0)\), as a program, runs in at most a constant
times
\[
  \frac{N^{2}(\log D+1)^{2}}{D^{\gamma}}+w\,D^{q(\theta)}(\log D+1)
\]
steps on every instance with \(1\le D\le N^{\varepsilon}\) and \(w\) wanted
positions.
\end{theorem}

The program is \lean{program31} of \cite{upstream}, unchanged.  Its time is
bounded upstream by a constant times any pair of bounds that dominate the two
costs of Theorem~30; the proof is repeated against \lean{cost8X}
(\lean{exists_tPreCore_leX}, \lean{exists_tOffline32_leX}), and the corollary
above supplies the domination (\lean{first_termX}, \lean{costsX}).  The theorem
is \lean{offlineWithinX}.  The statement that a procedure of the program solves
the thin product within these bounds on all instances with \(D\le N^{\varepsilon}\)
is \lean{thinClaim_gammaX}; it has two more hypotheses, which come from the
regime test \(D^{18}\le N\) in the program text and are discussed in
Section~\ref{sec:consequences}.

\subsection{The form of Section~2}

The same count improves Lemma~11 of the paper directly.  The theorems of
Section~\ref{sec:consequences} use the form of Section~4 above and do not need
the following statement.

\begin{proposition}[Lemma~11 with the exact count]\label{prop:lemmaA}
Let \(c\ge1\), \(\kappa\in\R\), \(m\ge1\) and \(L\ge\max(10m,cm)\).  Let \(W\)
be a set of wanted positions of an \(N\times N\) matrix with
\(|W|\le N^{2}/D^{\kappa}\), and for each tile \(T\) let \(W_T\) be the set of
output strings of \(T\) that are wanted.  Then
\[
  \sum_T|\Leaves(W_T)|\le\sum_{d=0}^{m}\min\bigl(|W|\alpha_d,\ \tau\beta_d\bigr)
  \le4(m+1)(L+1)\,N'^{2}\,e^{m\,\Gamma_c(\kappa)},
\]
where \(\tau\) is the number of tiles, \(N'\) is the size of the matrix padded
to a multiple of the side of a tile, and
\[
  \Gamma_c(\kappa)=\sup_{0\le\delta\le1}
    \min\bigl(\fone(\delta),\gtwo(c,\delta)\bigr),\qquad
  \fone(\delta)=-\kappa\ln 4+\Hent(\delta)+\delta\ln 9=(q(\delta)-\kappa)\ln 4 .
\]
If the side of a tile is at most \(N\), then \(N'\le2N\).
\end{proposition}

So the number of leaves is at most \(16(m+1)(L+1)N^{2}D^{-\gamma_c(\kappa)}\) with
\(\gamma_c(\kappa)=-\Gamma_c(\kappa)/\ln 4\).

\begin{proof}
By \eqref{eq:lemma11}, \(|\Leaves(W_T)|\le\sum_d\min(|W_T|\alpha_d,\beta_d)\).
A sum of minima is at most the minimum of the sums, and
\(\sum_T|W_T|\le|W|\); this is the first inequality.  Fix \(d\) and put
\(\delta=d/m\).  First, \(\alpha_d=\binomial md9^{d}\le
e^{m(\Hent(\delta)+\delta\ln 9)}\), so
\(|W|\alpha_d\le N^{2}e^{m\fone(\delta)}\le N'^{2}e^{m\fone(\delta)}\).
Second, \(\tau M\le4N'^{2}\), and by the inequality \eqref{eq:beta-int} with
\(t=d\) and by \eqref{eq:exp-g2},
\(\beta_d/M\le(L+1)e^{m\,\gtwo(L/m,\delta)}\le(L+1)e^{m\,\gtwo(c,\delta)}\), where
the last step is Lemma~\ref{lem:g2}(ii).  So each of the \(m+1\) values of \(d\)
contributes at most
\(4(L+1)N'^{2}e^{m\min(\fone(\delta),\gtwo(c,\delta))}\).
\end{proof}

In Lean (file \lean{ExactCount/Leaves.lean}): \(\fone\) is \lean{f1} and
\(\Gamma_c(\kappa)\) is \lean{GammaMM}; the first inequality is
\lean{sum_card_Leaves_le_sum_min}, the second is
\lean{sum_card_Leaves_le_exp_padN}, and the form with \(16N^{2}\) is
\lean{sum_card_Leaves_le_exp}.

\begin{remark}\label{rem:same-exponent}
\(\fone\) increases strictly on \([0,0.9]\) and \(\gtwo(c,\cdot)\) is
nonincreasing, so the supremum is taken where they cross: for every
\(\delta_0\in[0,0.9]\),
\(\min(\fone,\gtwo)(\delta_0)\le\Gamma_c(\kappa)\le\max(\fone,\gtwo)(\delta_0)\)
(\lean{min_le_GammaMM}, \lean{GammaMM_le_max}).  For \(c>10\) and
\(0<\kappa<\log_4 10\) they do cross, and \(\Gamma_c(\kappa)<0\)
(\lean{exists_crossing_f1_g2}).  At the crossing
\(\gX(c,\delta)=\kappa-q(\delta)\), which is also where the two terms of
\eqref{eq:cor32} balance when \(\gamma\) is replaced by \(\gX\).  The two forms
have the same exponent.
\end{remark}

\section{Pruned encodings}
\label{sec:pruned}

The second summand of \eqref{eq:cost8} pays for the encodings of the input at
all \(10^{L}\) leaves.  Most of them are never read.  This section computes
only those that are.  First the mathematics: the correctness of the pruned
recursion, its number of operations, and the resulting cost of the method
(Theorem~\ref{thm:pruned}, Proposition~\ref{prop:costP}).  Then the program:
the pruned encoder is written in the language of \cite{upstream}, the
procedures of the paper's data structure are re-verified under the weaker
contract that the encodings are right only at the leaves that are read, and
the running time of the new solver is proved on the machine
(Theorem~\ref{thm:encoder}).  Everything is machine-checked.

\subsection{What is read and what is nonzero}

Let \(a\) be the array of a band.  It is indexed by the \emph{left strings}:
strings \(u=u_1\cdots u_L\) of left variables.  For a leaf
\(\tau=\tau_1\cdots\tau_L\) the encoding is
\[
  \Phi_\tau(a)=\sum_u a[u]\prod_{\ell=1}^{L}\varphi_{\tau_\ell}(u_\ell),
\]
where \(\varphi_\tau\) is the linear form of the term \(\tau\) in the left
variables.  Each \(\varphi_\tau\) has at most three nonzero coefficients.  Two
facts hold in the recursion.
\begin{itemize}
\item Only leaves with at most \(m\) symbols \(\Pz\) are read: the leaves of
  \(\Leaves(U)\) have order at least \(0\).
\item \(a[u]=0\) unless \(u\) has exactly \(m\) inner variables.
\end{itemize}

Yates' algorithm computes the encodings level by level.  Stage \(k\) holds
\[
  A_k[\tau_1\cdots\tau_k;u_{k+1}\cdots u_L]
  =\sum_{u_1,\dots,u_k}a[u]\prod_{\ell\le k}\varphi_{\tau_\ell}(u_\ell),
\]
so \(A_0=a\) and \(A_L[\tau]=\Phi_\tau(a)\), and
\(A_{k+1}[\tau_1\cdots\tau_{k+1};\cdot]
=\sum_v\varphi_{\tau_{k+1}}(v)A_k[\tau_1\cdots\tau_k;v\,\cdot]\).

\begin{definition}\label{def:pruned}
Let \(P_k\) be the set of strings of \(k\) terms with at most \(m\) symbols
\(\Pz\), and \(U_k\) the set of strings of \(L-k\) left variables with at most
\(m\) inner variables.  The \emph{pruned recursion} computes at stage \(k\)
only the entries of \(A_k\) with index in \(P_k\times U_k\), by the rule above,
and treats all other entries as \(0\).
\end{definition}

\subsection{Correctness and cost}

\begin{theorem}[Pruned encodings]\label{thm:pruned}
Let \(L\ge10m\), and let \(a\) vanish off the left strings with exactly \(m\)
inner variables.
\begin{enumerate}
\item The pruned recursion computes \(\Phi_\tau(a)\) for every leaf \(\tau\)
  with at most \(m\) symbols \(\Pz\).
\item \(|P_k|=\sum_{j\le m}\binomial kj9^{k-j}\) and
  \(|U_k|=\sum_{i\le m}\binomial{L-k}{i}4^{i}3^{L-k-i}\).
\item \(\sum_{k=0}^{L}|P_k|\,|U_k|\le\frac92|P_L|\le\frac92\cdot\frac{M}{1-\rho}\)
  with \(\rho=9m/(L-m+1)\).
\item Every computed entry is a sum of at most three products.  So the pruned
  recursion makes at most \(\frac{27}{2}M/(1-\rho)\) multiplications.
\end{enumerate}
\end{theorem}

\begin{proof}
(i) An entry of \(A_k\) whose suffix has more than \(m\) inner variables is
\(0\): every \(u\) in its sum has more than \(m\) inner variables, so
\(a[u]=0\).  An entry of \(A_{k+1}\) with prefix in \(P_{k+1}\) reads entries of
\(A_k\) whose prefix is in \(P_k\), since deleting a term does not add a
symbol \(\Pz\).  So, by induction on \(k\), the pruned arrays agree with those
of Yates' algorithm at every index whose prefix has at most \(m\) symbols
\(\Pz\): such an entry is \(0\) if its suffix is outside \(U_k\), and otherwise
each entry it reads either lies in \(P_{k-1}\times U_{k-1}\) or is \(0\) in both
arrays.

(ii) There are nine terms other than \(\Pz\), and four inner and three outer
left variables.

(iii) Since \(\binomial{k+1}{j}\ge\binomial kj\), we have
\(|P_{k+1}|\ge9|P_k|\).  By Pascal's rule,
\[
  \sum_{i\le m}\binomial{n+1}{i}4^{i}3^{n+1-i}
  =3\sum_{i\le m}\binomial ni4^{i}3^{n-i}
   +4\sum_{i\le m-1}\binomial ni4^{i}3^{n-i},
\]
so \(|U_k|\le7|U_{k+1}|\).  With \(|U_L|=1\) this gives
\(|P_k||U_k|\le(7/9)^{L-k}|P_L|\), and \(\sum_{k}(7/9)^{L-k}\le9/2\).  The terms
of \(|P_L|=\sum_{j\le m}\binomial Lj9^{L-j}\) have consecutive ratios
\(\binomial L{j-1}9/\binomial Lj=9j/(L-j+1)\le\rho<1\), and the term with
\(j=m\) is \(M\).  So \(|P_L|\le M/(1-\rho)\).

(iv) follows from (iii) and the three coefficients of a form.
\end{proof}

In Lean (directory \lean{PrunedEncoding}): Yates' algorithm is
\lean{yatesStage} with \lean{yatesStage_final}; the pruned recursion is
\lean{prunedStage}; part~(i) is \lean{prunedStage_eq_yatesStage} and
\lean{prunedStage_final}; part~(ii) is \lean{card_stageSet_of_le} with
\lean{prefCount} and \lean{sufCount}; part~(iii) is
\lean{sum_card_stageSet_le}; part~(iv) is \lean{prunedMults_le_M}; the theorem
as a whole is \lean{pruned_encoding}.

\begin{remark}\label{rem:PL}
\(|P_L|=\sum_{d=0}^{m}\beta_d\): the pruned recursion computes exactly the
encodings at the leaves of order at least \(0\).  For \(L\ge cm\) with \(c>10\)
fixed, \(1/(1-\rho)\le(c-1)/(c-10)\).
\end{remark}

\begin{remark}[The work of the program]\label{rem:encwork}
The program of Section~\ref{sec:encoder-program} below follows the recursion
of Theorem~\ref{thm:pruned} with one simplification: at stage \(k\) it keeps
the entries with index in \(P_k\times(\text{all \(7^{L-k}\) suffixes})\)
instead of \(P_k\times U_k\), so that the suffixes are indexed densely and no
compact indexing of \(U_k\) is needed.  For \(L\ge10m\) the largest slice,
with \(7^{L}\) entries, is no larger than a tile: \(7^{L}\le9^{L-m}\le M\)
(\lean{seven_pow_le_M}).  The number of computed entries, \lean{encWorkP}~\(L\,m\)
in Lean, is
\[
  \sum_{k=0}^{L}|P_k|\,7^{L-k}\le\frac92|P_L|
  \le\frac92\cdot\frac{M}{1-\rho},
\]
from \(9|P_k|\le|P_{k+1}|\) and \(\sum_i(7/9)^i\le9/2\)
(\lean{encWorkP_le}, \lean{encWorkP_le_M_div}): part~(iii) with \(7^{L-k}\)
in place of \(|U_k|\), and the same bound.  The time of the program is at most
\(\EncConst\) times this number (\lean{tEncP_le}).
\end{remark}

\subsection{The thinness bound}

With \(O(L)\) word operations per entry for the index arithmetic, a band costs
\(O(LM)\) instead of \(O(L\,10^{L})\), and all \(N/(\sqrt K N_0)\) bands cost
\(O(L\,N\sqrt K N_0)\), since \(M=KN_0^{2}\).  This is at most
\(N^{2}/D^{\gamma}\) up to the factor \(L\) when \(\sqrt K N_0D^{\gamma}\le N\).
So the thinness condition \eqref{eq:Rc} becomes
\begin{equation}\label{eq:RcP}
  \varepsilon<\RcP(\gamma):=\frac{\ln 4}{\Apr(c)+\gamma\ln 4},\qquad
  \Apr(c)=\frac c2\Hent\Bigl(\frac1c\Bigr)+(c-1)\ln 3 ,
\end{equation}
where \(\Apr(c)\) is the exponential rate of \(\sqrt K N_0\) per inner level.
Since \(M\le10^{L}\), we have \(\Apr(c)\le\Afull(c)\) and
\(\RcP(\gamma)\ge\Rc(\gamma)\).  In Lean the constants are \lean{basePruned} and
\lean{baseFull}, the bound is \lean{RcP}, and the two inequalities are
\lean{basePruned_le_baseFull} and \lean{Rc_le_RcP}.

\begin{proposition}[The costs with pruned encodings]\label{prop:costP}
Let \(c>10\), \(0<\theta<0.9\), \(0\le\gamma<\gX(c,\theta)\) and
\(0<\varepsilon<\RcP(\gamma)\).  There are \(C\) and \(m_0\) such that the
following holds for all \(D\ge2\) and \(N\) with \(D\le N^{\varepsilon}\) and
\(m=\lceil\log_4 D\rceil\ge m_0\), with \(L=\lceil cm\rceil\) and
\(t=\lceil\theta m\rceil\).
\begin{enumerate}
\item The tile fits: \(\sqrt K N_0\le N\).
\item The cost of the preprocessing with pruned encodings,
  \[
    L\,m\,\frac{\sum_{d=t}^{m}\beta_d}{M}\,N^{2}
    +\frac{N}{\sqrt K N_0}\cdot L\cdot\frac{27}{2}\cdot\frac{M}{1-\rho},
  \]
  is at most \(C\,N^{2}\log^{2}D/D^{\gamma}\).
\item The cost \(L\sum_{d\le t}\alpha_d\) of a query is at most
  \(C\,D^{q(\theta)}\log D\).
\end{enumerate}
\end{proposition}

\begin{proof}
The first summand of (ii) is bounded as in Section~\ref{sec:exact}.  The second
is \(\frac{27}{2}NL\sqrt K N_0/(1-\rho)\).  Now
\(\sqrt K N_0\le e^{m\Apr(L/m)}\), the ratio \(L/m\) tends to \(c\), and
\(N\ge D^{1/\varepsilon}>4^{(m-1)/\varepsilon}\); so
\(D^{\gamma}L\sqrt K N_0\le N\) for all large \(m\), because
\(\Apr(c)+\gamma\ln 4<\ln 4/\varepsilon\).  This gives (i) and (ii).  Part~(iii)
is as in the paper.
\end{proof}

In Lean (file \lean{PrunedEncoding/Costs.lean}): the expression of (ii) is
\lean{cost8P}, the inequality \(D^{\gamma}L\sqrt K N_0\le N\) is \lean{eq_10P},
and the proposition is \lean{costsP}.  It is the analogue of the statement
\lean{costsX} on which Theorem~\ref{thm:offline} rests, with the pruned
encoder's count of operations in place of \(10^{L}\) and \(\RcP\) in place of
\(\Rc\).

\begin{theorem}[The encoder program]\label{thm:encoder}
For all rational \(c=a/b>10\) and \(\theta=p/q\in(0,0.9)\), every
\(0<\gamma<\gX(c,\theta)\) and every \(\varepsilon<\RcP(\gamma)\), there is a
procedure, in the programming language of \cite{upstream}, that solves the thin
product on all instances and takes at most a constant times
\[
  \frac{N^{2}(\log D+1)^{2}}{D^{\gamma}}+w\,D^{q(\theta)}(\log D+1)
\]
steps on every instance with \(1\le D\le N^{\varepsilon}\) and at most \(w\)
wanted positions.
\end{theorem}

Theorem~\ref{thm:encoder} is the conclusion of Theorem~\ref{thm:offline} with
\(\RcP\) in place of \(\Rc\).  On paper it follows from
Theorem~\ref{thm:pruned} and Proposition~\ref{prop:costP}: replace the encoder
of the offline routine by the pruned recursion, and replace the regime test of
the program, which is tied to \(\varepsilon\le1/18\), by the test with a
rational exponent of Section~\ref{sec:regime}.  In Lean the statement is the
proposition \lean{PrunedEncoderClaim} (file \lean{Pipeline/General.lean}),
proved as \lean{prunedEncoderClaim} (file \lean{Pipeline/PrunedClaim.lean}).
It is stated with the same notion, \lean{ThinClaim}, in which
Theorems~\ref{thm:T3} and~\ref{thm:T4} obtain the thin product from the
existing program (\lean{thinClaim_gammaX}, \lean{thinClaimRS_gammaX}): a claim
about the time model
\lean{lightModel} of the upstream programs, which \cite{upstream} compiles to
the word RAM.  The rest of this section describes the program and its proof.

\subsection{The encoder program}
\label{sec:encoder-program}

\emph{The encoder.}  Upstream's encoder is a recursive procedure
\lean{encodeBody} that descends the ten terms of the identity level by level and
fills the \(10^{L}\) encodings of a band.  The pruned encoder
(\lean{encodePBody}, file \lean{PrunedProgram/Encode/Recursion.lean}) is the
same procedure with a budget \(j\) of symbols \(\Pz\): a loop over the nine
terms \(\lambda\ne\Pz\) recurses with the same budget, and the term \(\Pz\)
(digit~\(9\)) is treated after the loop, only if the budget is nonzero, with
budget \(j-1\).  The last term is treated outside the loop because the
language's \texttt{for} charges every round the same time.  The slices are the
dense ones of Remark~\ref{rem:encwork}.  Its correctness is \lean{encodeP_spec}:
after the call with budget \(j\), every leaf with at most \(j\) symbols \(\Pz\)
holds its encoding (the postcondition \lean{EncodePostP}), and its time is at
most \(\EncConst\) times the work \lean{encWorkP}~\(L\,j\) (\lean{tEncP_le}),
which Remark~\ref{rem:encwork} bounds by
\(\frac92\cdot\EncConst\cdot M/(1-\rho)\) for \(j=m\).

\emph{The weaker contract.}  Upstream's data structure assumes that the
encodings of a band are right at all \(10^{L}\) leaves (a \emph{segment}).
With the pruned encoder they are right only at the leaves with at most \(m\)
symbols \(\Pz\); this is the predicate \lean{SegOn} (file
\lean{PrunedProgram/SegOn.lean}).  The preprocessing and the query of
Theorem~30 of the paper read an encoding cell in exactly two places, both at
such leaves: in the construction of the trie, at the leaf of a code of
\lean{nineStrs}, which has at most \(m-t\) digits~\(9\); and in a query, at a
leaf that contributes to the output string, which has the symbol \(\Pz\) only
at inner levels (\lean{P0Levels_subset_innerSetO},
\lean{card_P0Levels_le_of_contributes}).  So upstream's own procedures for the
tiles (\lean{fillListBody}, \lean{tileBody}, \lean{allTilesBody}) and for a query
(\lean{queryAtBody}) are re-verified under the weaker contract, with their program
text unchanged: \lean{fillListP_meets} and \lean{preCoreP_of_base58} (files
\lean{PrunedProgram/Tiles}, \lean{PrunedProgram/Pre/Preprocessing.lean}),
\lean{queryCoreP_spec} and \lean{queryAtP_of_base58} (directory
\lean{PrunedProgram/Query}).  Upstream's shared stage, which builds the
tables and encodes every band, becomes \lean{sharedPBody} with the entry
\lean{sharedP_entry} (file \lean{PrunedProgram/Shared/Stage.lean}): its two
calls of the encoder are exchanged for calls of the pruned one, and nothing
else changes.  The invariant of the finished data structure is \lean{DSReadyP}
(file \lean{PrunedProgram/Pre/Defs.lean}).

\emph{The program.}  The routines with rational parameters are
\lean{pre31PBody}, \lean{offline32PBody} and \lean{offline32P_meets}
(directory \lean{PrunedProgram/Wrapper}).  The program is \lean{programP}: the
program \lean{programRS} of Section~\ref{sec:regime} followed by seven
procedures, \lean{encodeP}, \lean{encodeBandsP}, \lean{sharedP},
\lean{preCoreP}, \lean{pre31P}, \lean{offline32P} and \lean{allP}; every other
procedure is upstream's, with its number.  The offline routine as an object of
the pipeline is \lean{offline32PRoutine}.

\emph{The cost.}  The running time is analyzed against \lean{cost8P} as the
running time of upstream's program is analyzed against \lean{cost8X} in
Theorem~\ref{thm:offline}: \lean{Within8P} and \lean{exists_tPreCore_leP},
\lean{CostsWithinP} and \lean{offlineWithinP} (directory \lean{Cost8P}) are
the analogues of \lean{exists_tPreCore_leX}, \lean{costsX} and
\lean{offlineWithinX}.  The small summands of the shared stage are bounded by
\(NL\sqrt K N_0\) directly, since \lean{cost8P} has no term \(10^{L}\).

\emph{The pipeline.}  The solver of all instances is generic in the offline
routine: the structure \lean{OfflineRoutine} (file
\lean{Pipeline/ThinClaim.lean}) packages the procedure, its time, what it
needs from the program and its specification, as the regime test is already
packaged as \lean{RegimeTest}.  Upstream's routine is the instance
\lean{offline32Routine} and the pruned one is \lean{offline32PRoutine};
\lean{thinClaimRS_of_offlineWithinT} (file \lean{Pipeline/RegimeRS.lean})
turns either into a thin-product claim at the regime \(D^{r}\le N^{s}\).  For
\lean{thinClaimP_gammaX} and \lean{prunedEncoderClaim} the fraction \(r/s\) is
chosen strictly between \(\Apr(c)/\ln 4\) and \(1/\varepsilon\), which is
possible because \(\varepsilon<\RcP(\gamma)<\ln 4/\Apr(c)\); the case
\(\varepsilon\le0\) of the statement is reduced to a positive \(\varepsilon\)
by \lean{ThinClaim.mono}.

\section{Charging the solver at the true size of the middle part}
\label{sec:middle}

In the reduction of Section~\ref{sec:host} the middle part of an instance is
\(C_k\times\Z_p\) with \(|C_k|\le\lceil s/g'\rceil\) and \(p\le s\),
\(s=\lfloor\sqrt{D'}\rfloor\).  So it has at most
\[
  \operatorname{mid}(D',g'):=\Bigl\lceil\frac{s}{g'}\Bigr\rceil s\approx\frac{D'}{g'}
\]
vertices.  The paper hands the instance to the solver as an instance of
\(\LopTri(n,D')\) and charges it accordingly.  Handing it over as an instance
of \(\LopTri(n,\operatorname{mid}(D',g'))\) changes one assignment of the
program and separates two roles of \(D'\): the number of query pairs is
governed by \(\sqrt{D'}\), the cost of a call by \(D'/g'\).

\subsection{The reduction}

\begin{theorem}[Theorem~17 at the true middle size]\label{thm:mid}
There is a constant \(C\) with the following property.  Let \(T(n,D,w)\) bound
the time of a solver of \(\LopTri(n,D)\) with \(w\) query pairs.  Then Exact
Triangle on \(n\) vertices per part with weights of absolute value at most
\(n^{\nu}\), \(\nu\ge1\), is solved in time at most
\[
  4ng'\Bigl(T\bigl(n,\operatorname{mid}(D',g'),\lfloor n^{2}/\sqrt{D'}\rfloor\bigr)
     +C\frac{n^{2}}{\sqrt{D'}}\Bigr)
  +C\Bigl(\frac{\nu n^{3}\log n}{g'}+n^{\log_2 7}D'^{3/2}+n^{2}D'g'\Bigr)
\]
whenever \(16\le D'\le n\) and \(1\le g'\le\sqrt{D'}\).
\end{theorem}

This is a statement about programs: the solver is a procedure of a program of
the machine, and the theorem gives a program for Exact Triangle.  The exponent
\(\log_2 7\) is that of Strassen's algorithm \cite{Str69}, which the upstream
program uses for the choice of the prime; see Section~\ref{sec:limits}.

In Lean the statement is \lean{Claim17Mid}, the upstream statement
\lean{Claim.Theorem_17} with \(\operatorname{mid}(D',g')\) (\lean{midSize}) in
place of \(D'\) in the first argument of \(T\).  The host is upstream's with one
added assignment (\lean{et17Sizes'}): after the prime and the size of a piece
are computed, the parameter that is passed to the solver is overwritten with
\(\lceil s/g'\rceil s\).  Its correctness and its time follow from upstream's by
rewriting the arguments of the solver (\lean{hostTime'_eq},
\lean{obeysBound17Mid_hostTime}); the claim for programs is
\lean{claim17Mid_of_paramProcs}, and for \(D'=\lfloor n^{a/b}\rfloor\),
\(g'=\lceil D'^{c/d}\rceil\) with rational exponents it is
\lean{hostBound_mid_rat}.

\subsection{The total time}

\begin{theorem}[Total time]\label{thm:total}
Let \(\gamma,q\ge0\) and \(\varepsilon>0\).  Suppose that \(\LopTri(n,D)\) with
\(w\) query pairs is solved, for \(D\le n^{\varepsilon}\), in time
\(O\bigl(n^{2}(\log D+1)^{2}/D^{\gamma}+w\,D^{q}(\log D+1)\bigr)\).  Let
\(0<d\) and \(0<\eta<1/2\) be rational with \(d(1-\eta)<\varepsilon\), and let
\(D'=\lfloor n^{d}\rfloor\), \(g'=\lceil D'^{\eta}\rceil\).  Then Exact Triangle
with weights up to \(n^{\nu}\) is solved in time \(O(\nu\,n^{3-\delta}\log n)\)
for every \(\delta\ge0\) with
\begin{equation}\label{eq:total}
  \delta<d\,\min\Bigl(\eta,\ \min\bigl(\gamma(1-\eta),\ \tfrac12-q(1-\eta)\bigr)-\eta\Bigr)
  \qquad\text{and}\qquad \tfrac32d+\delta\le0.19 .
\end{equation}
\end{theorem}

\begin{proof}
Apply Theorem~\ref{thm:mid}.  The dimension of a call is
\(D=\operatorname{mid}(D',g')=\Theta(D'^{1-\eta})=\Theta(n^{d(1-\eta)})\), which
is at most \(n^{\varepsilon}\) for large \(n\).  Up to constants and powers of
\(\log n\), the terms are
\begin{align*}
  &\text{calls, preprocessing:} & n\,g'\,n^{2}D^{-\gamma}&=n^{3+d\eta-d\gamma(1-\eta)},\\
  &\text{calls, queries:} & n\,g'\,\frac{n^{2}}{\sqrt{D'}}D^{q}&=n^{3+d\eta+dq(1-\eta)-d/2},\\
  &\text{scans:} & \nu n^{3}/g'&=\nu n^{3-d\eta},\\
  &\text{prime:} & n^{\log_2 7}D'^{3/2}&\le n^{2.81+3d/2},\\
  &\text{construction:} & n^{2}D'g'&=n^{2+d+d\eta}.
\end{align*}
The first three exponents are at most \(3-\delta\) by the first condition of
\eqref{eq:total}, with room for the logarithms, and so is that of the term
\(ng'n^{2}/\sqrt{D'}\).  The last two are at most \(3-\delta\) by the second
condition.
\end{proof}

In Lean: \lean{explicitFrom_mid}, from \lean{hostBound_mid_rat} and a solver
of \(\LopTri\) (\lean{LopClaim}, obtained from a thin-product claim by
\lean{lopClaim_of_thinClaim}).  The same calculation with every call charged at
\(D'\), as in the paper, is \lean{explicitFrom_original}; its condition is
\(\delta<d\min(\eta,\min(\gamma,\frac12-q)-\eta)\).

\subsection{The saving of a tuple}

It is convenient to take the dimension \(D\) of the calls as the basic
parameter.  Write \(D\approx n^{\varepsilon}\) and
\begin{equation}\label{eq:sigma}
  \sigma=\frac{1}{2(1-\eta)}\in\bigl[\tfrac12,1\bigr),\qquad\text{so that}\qquad
  D'=D^{2\sigma},\quad g'=D^{2\sigma-1},\quad\sqrt{D'}=D^{\sigma}.
\end{equation}
An instance then has inner dimension \(D\) and at most \(n^{2}/D^{\sigma}\)
query pairs.  With the thin product of Section~\ref{sec:exact} at the
parameters \((c,\theta)\) the exponent of a call is
\(\min(\gX(c,\theta),\sigma-q(\theta))\), by \eqref{eq:cor32} with
\(\kappa=\sigma\), and the first condition of \eqref{eq:total} reads
\(\delta<\sav(c,\theta,\sigma,\varepsilon)\) with \(d(1-\eta)\) in the place of
\(\varepsilon\), where:

\begin{definition}\label{def:saving}
The \emph{saving} of the tuple \((c,\theta,\sigma,\varepsilon)\) is
\[
  \sav(c,\theta,\sigma,\varepsilon)
  =\varepsilon\,\min\Bigl(2\sigma-1,\ \min\bigl(\gX(c,\theta),\sigma-q(\theta)\bigr)-(2\sigma-1)\Bigr).
\]
\end{definition}

In Lean: \lean{tupleSaving}, with \lean{effGamma} for the inner minimum.  The
first argument of the outer minimum is the saving of the scans, the second
that of the calls.  In the paper every call is charged at the dimension \(D'\); relative to
\(D'\) the density of the query pairs is \(\kappa=1/2\).  Relative to the true
dimension \(D\) it is \(\sigma>1/2\) as soon as \(g'>1\), and the thin product
profits from both: a smaller inner dimension and, relative to it, fewer wanted
entries.

\begin{remark}[A larger hash prime]\label{rem:prime}
The same algorithm can be described without \(g'\).  Hash modulo a prime
\(p\approx D^{\sigma}\) with \(\sigma\ge1/2\), and cut each part into pieces of
\(D^{1-\sigma}\) vertices, so that a middle part \(C_k\times\Z_p\) has \(D\)
vertices.  There are \(O(nD^{2\sigma-1})\) instances with at most
\(n^{2}/D^{\sigma}\) query pairs each, and the scans cost
\(\nu n^{3}\log n/D^{2\sigma-1}\).  This is the reduction above with
\(D'=D^{2\sigma}\) and \(g'=D^{2\sigma-1}\): a larger prime and no further
splitting.  The balance of Remark~20 of the paper, \(\eta=\gamma/2\), becomes
the fixed point \(\sigma=1/2+\gamma/4\) of Section~\ref{sec:optimum}.
\end{remark}

\section{The optimal parameters}
\label{sec:optimum}

The method has four real parameters: the ratio \(c>10\), the switching ratio
\(0<\theta<0.9\), the size \(1/2\le\sigma<1\) of the prime, and the exponent
\(\varepsilon>0\) of the inner dimension.  This section finds the supremum of
the saving of Definition~\ref{def:saving} over the parameters that the method
allows.

\subsection{Definitions}

Let \(A>0\) be a \emph{thinness constant} and
\[
  R(\gamma)=\frac{\ln 4}{A+\gamma\ln 4},\qquad F(\gamma)=\frac{R(\gamma)\,\gamma}{2}.
\]
The paper's encodings have \(A=\Afull(c)\), so that \(R=\Rc\) of \eqref{eq:Rc},
and pruned encodings have \(A=\Apr(c)\), so that \(R=\RcP\) of \eqref{eq:RcP}.
A tuple \((c,\theta,\sigma,\varepsilon)\) is \emph{admissible} for \(A\) if
\(c>10\), \(0<\theta<0.9\), \(1/2\le\sigma<1\) and
\(0<\varepsilon<R(\gX(c,\theta))\).  Put
\[
  \gQ(\theta)=\frac{2-4q(\theta)}{3},\qquad
  \Gam(c)=\sup_{0<\theta<0.9}\min\bigl(\gX(c,\theta),\gQ(\theta)\bigr).
\]
In Lean: \lean{thinR}, \lean{savingF}, \lean{AdmissibleX}, \lean{gammaQ},
\lean{Gam}.

\subsection{The crossing point}

\begin{lemma}\label{lem:crossing}
Let \(c>10\).  On \([0,0.9]\) the function \(\gX(c,\cdot)\) is continuous and
strictly increasing with \(\gX(c,0)=0\), and \(\gQ\) is continuous and strictly
decreasing with \(\gQ(0)=2/3\) and \(\gQ(0.9)<0\).  They cross at exactly one
point \(\theta^{*}(c)\in(0,0.9)\), and \(\Gam(c)\) is their common value there.
Moreover \(0<\Gam(c)<2/3\), and
\[
  \min\bigl(\gX(c,\theta),\gQ(\theta)\bigr)\le\Gam(c)
  \le\max\bigl(\gX(c,\theta),\gQ(\theta)\bigr)\qquad(0<\theta<0.9).
\]
\(\Gam\) is nondecreasing on \((10,\infty)\).
\end{lemma}

\begin{proof}
\(\gX(c,\cdot)\) increases strictly by Lemma~\ref{lem:g2}(i).  The function
\(q\) increases strictly on \([0,0.9]\), since
\(q'(\theta)\ln 4=\ln(9(1-\theta)/\theta)>0\) for \(\theta<0.9\), and
\(q(0.9)=\log_4 10>1\).  The intermediate value theorem gives the crossing, and
monotonicity gives uniqueness and the two inequalities.  The last statement
holds because \(\gX\) is nondecreasing in \(c\) (Lemma~\ref{lem:g2}(ii)), so
that \(\min(\gX(c,\theta),\gQ(\theta))\) is.
\end{proof}

In Lean: \lean{existsUnique_crossing}, \lean{Gam_eq_of_crossing},
\lean{Gam_pos}, \lean{Gam_lt_two_thirds}, \lean{Gam_ge_min}, \lean{Gam_le_max},
\lean{Gam_monotoneOn}.  The two inequalities are what the numerical bounds
below rest on: every \(\theta\) gives a lower and an upper bound for
\(\Gam(c)\).

\subsection{The supremum for a fixed ratio}

\begin{theorem}[Supremum of the saving]\label{thm:sup}
Let \(c>10\) and \(A>0\).  Over the tuples \((c,\theta,\sigma,\varepsilon)\)
that are admissible for \(A\),
\[
  \sup\ \sav(c,\theta,\sigma,\varepsilon)=S(c):=F(\Gam(c))=\frac{R(\Gam(c))\,\Gam(c)}{2},
\]
and the supremum is not attained.  It is approached at
\(\theta=\theta^{*}(c)\), \(\sigma=1/2+\Gam(c)/4\) and
\(\varepsilon\uparrow R(\Gam(c))\).
\end{theorem}

\begin{proof}
\(R\) is positive and strictly decreasing on \([0,\infty)\), and
\(F(\gamma)=\frac12\bigl(1-A/(A+\gamma\ln 4)\bigr)\) is strictly increasing
there.

\emph{Upper bound.}  Let the tuple be admissible, and put
\(a=\gX(c,\theta)\), \(u=2\sigma-1\in[0,1)\) and
\(v=\min(u,\min(a,\sigma-q(\theta))-u)\), so that the saving is
\(\varepsilon v\).  If \(v\le0\) the saving is at most \(0<F(\Gam(c))\).  Let
\(v>0\).  Then \(v\le u\) and \(v\le a-u\), so \(v\le a/2\).  Also
\(v\le\sigma-q(\theta)-u=(1-u)/2-q(\theta)\le(1-v)/2-q(\theta)\), so
\(v\le(1-2q(\theta))/3=\gQ(\theta)/2\).  Hence
\(v\le\min(a,\gQ(\theta))/2\le\Gam(c)/2\) by Lemma~\ref{lem:crossing}.  Since
\(\varepsilon<R(a)\) and \(v>0\), the saving is less than \(R(a)v\).  If
\(a\le\Gam(c)\), then \(R(a)v\le R(a)a/2=F(a)\le F(\Gam(c))\).  If
\(a>\Gam(c)\), then \(R(a)v\le R(\Gam(c))\Gam(c)/2\).

\emph{Approach.}  At \(\theta=\theta^{*}(c)\) we have
\(\gX(c,\theta)=\gQ(\theta)=\Gam\), that is, \(q(\theta)=1/2-3\Gam/4\).  With
\(\sigma=1/2+\Gam/4\) we get \(u=\Gam/2\) and \(\sigma-q(\theta)=\Gam\), so
\(v=\Gam/2\) and the saving is \(\varepsilon\Gam/2\), for every
\(0<\varepsilon<R(\Gam)\).  Here \(\sigma<1\) because \(\Gam<2/3\).
\end{proof}

In Lean: \lean{tupleSaving_lt_savingF}, \lean{tupleSaving_crossing},
\lean{exists_tupleSaving_gt}, \lean{isLUB_tupleSaving} and
\lean{savingF_notMem_tupleSaving}.

Programs can only use rational parameters.  Nothing is lost: the saving is
continuous in the four parameters, admissibility is an open condition when
\(A\) depends continuously on \(c\), and both \(\Afull\) and \(\Apr\) do.

\begin{proposition}\label{prop:rational}
Let \(A\) be \(\Afull\) or \(\Apr\), let \(c>10\) and \(s<F_{A(c)}(\Gam(c))\).
For every \(c_1>c\) there are rational numbers
\(c',\theta',\sigma',\varepsilon'\) with \(c<c'<c_1\) and \(\sigma'>1/2\) such
that \((c',\theta',\sigma',\varepsilon')\) is admissible for \(A(c')\) and
\(\sav(c',\theta',\sigma',\varepsilon')>s\).
\end{proposition}

In Lean: \lean{exists_rat_tuple_baseFull} and
\lean{exists_rat_tuple_basePruned}, from \lean{continuous_tupleSaving}.

If every call is charged at \(D'\), as in the paper, the condition of
Theorem~\ref{thm:total} is \(\delta<d\min(\eta,\min(\gamma,\frac12-q)-\eta)\)
with \(d<\varepsilon\).  The constraint from the queries is then
\(\gamma\le\frac12-q\) instead of \(\gamma\le\sigma-q\) with
\(\sigma=\frac12+\frac\gamma4\), which is \(\gamma\le\gQ\).  The following
proposition says what this costs.

\begin{proposition}[The paper's accounting]\label{prop:half}
Let \(c>10\) and \(A>0\).  Put
\(\Gam_{1/2}(c)=\sup_{0<\theta<0.9}\min\bigl(\gX(c,\theta),\tfrac12-q(\theta)\bigr)\),
and call \((\theta,\gamma,\varepsilon,\eta)\) admissible for \(A\) if
\(0<\theta<0.9\), \(0<\gamma<\gX(c,\theta)\), \(0<\varepsilon<R(\gamma)\) and
\(0<\eta<\frac12\).  Then \(\gX(c,\cdot)\) and \(\frac12-q\) cross in
\((0,0.9)\), \(\Gam_{1/2}(c)\) is their value at the crossing, and
\[
  \sup\ \varepsilon\min\bigl(\eta,\ \min(\gamma,\tfrac12-q(\theta))-\eta\bigr)
  =F(\Gam_{1/2}(c))<F(\Gam(c)),
\]
the supremum over the admissible tuples; moreover \(\Gam_{1/2}(c)<\Gam(c)\).
\end{proposition}

The admissible tuples are exactly the hypotheses of Theorem~\ref{thm:T3}'s
general form (\lean{allSolved_exact}), with \(d\) absorbed into
\(\varepsilon\).  The supremum is approached at the crossing with
\(\eta=\gamma/2\).  In Lean (file \lean{Optimum/PaperAccounting.lean}):
\lean{GamHalf}, \lean{exists_crossing_half}, \lean{GamHalf_eq_of_crossing},
\lean{GamHalf_lt_Gam}; the saving is \lean{paperSaving}, the admissible tuples
are \lean{AdmissibleHalf}, the supremum is \lean{isLUB_paperSaving}, and the
strict inequality between the two suprema is \lean{savingF_GamHalf_lt}.

\subsection{Numerical values and the optimum over the ratio}
\label{sec:numerics}

Both thinness constants are nondecreasing in \(c\) on \([10,\infty)\): the
derivatives are
\(\Afull'(c)=\ln\frac{10}{3}+\frac12\ln(1-\frac1c)\) and
\(\Apr'(c)=\ln 3-\frac12\ln(1-\frac1c)\) (\lean{baseFull_monotoneOn},
\lean{basePruned_monotoneOn}).  So two effects compete in \(S(c)=F_{A(c)}(\Gam(c))\): \(\Gam\) grows with
\(c\), and \(R\) shrinks.

Numerically, with pruned encodings \(S\) is largest at \(c\approx\cPrApprox\),
where \(S\approx\SprApprox\), \(\Gam\approx\GamPrApprox\),
\(\sigma\approx\sigmaPrApprox\), \(\varepsilon\approx\epsPrApprox\) and
\(\theta^{*}\approx\thetaPrApprox\).  With the paper's encodings it is largest
at \(c\approx\cFullApprox\), where \(S\approx\SfullApprox\).

These values are computed in floating point.  The following are certified.

\begin{proposition}[Values at three ratios]\label{prop:values}
\begin{enumerate}
\item At \(c=\cUncondOpt\) with the paper's encodings:
  \(24.8438\le\Afull(c)\le24.84381\), \(0.0741364\le\Gam(c)\le0.0741372\) and
  \(\SfullLo<S(c)<0.00206\).
\item At \(c=\cUncond\) with the paper's encodings: \(\Afull(c)\ge18\ln4\),
  \(0.074835\le\Gam(c)\le0.074837\), \(0.05504<\Rc(\Gam(c))<0.05505\) and
  \(\dFour<S(c)<0.00206\).
\item At \(c=\cCond\) with pruned encodings: \(0.0762<\Gam(c)<0.0763\),
  \(\epsCond<\RcP(\Gam(c))<0.05499\) and \(\SstarLo<S(c)<\SstarHi\).
\end{enumerate}
\end{proposition}

In Lean: \lean{baseFull_107_5_mem}, \lean{Gam_107_5_mem},
\lean{savingF_full_107_5_gt}, \lean{savingF_full_107_5_lt}; \lean{baseFull_ge_18},
\lean{Gam_108_5_mem}, \lean{Rc_108_5_bounds}, \lean{savingF_full_gt},
\lean{savingF_full_lt}; \lean{Gam_22_bounds}, \lean{thinR_pruned_bounds},
\lean{savingF_pruned_gt}, \lean{savingF_pruned_lt}.

\begin{theorem}[The optimum over the ratio]\label{thm:global}
Let \(S(c)=F_{\Apr(c)}(\Gam(c))\) be the supremum of the saving with pruned
encodings at the ratio \(c\).  Then
\[
  \SstarLo<\sup_{c>10}S(c)\le\SstarHi,
\]
and every \(c>10\) with \(S(c)\ge\SstarLo\) lies in \((\cStarLo,\cStarHi)\).
\end{theorem}

\begin{proof}
The lower bound is Proposition~\ref{prop:values}(iii).  For the upper bounds,
every \(\theta\in(0,0.9)\) gives \(\Gam(c)\le\max(\gX(c,\theta),\gQ(\theta))\)
by Lemma~\ref{lem:crossing}, and the values of \(\gX\), \(\gQ\) and \(\Apr\) at
rational arguments are enclosed between rational numbers by enclosing the
logarithms that occur in them.  On a cell \([c_1,c_2]\) both \(\Apr\) and
\(\Gam\) are nondecreasing, \(F_A(\gamma)\) decreases in \(A\) and increases in
\(\gamma\), so
\[
  S(c)\le F_{\Apr(c_1)}(\Gam(c_2))\qquad(c_1\le c\le c_2).
\]
The interval \((10,200]\) is cut into 222 cells, each with a certified bound of
this form: 68 cells cover \((10,\cStarLo]\) and \([\cStarHi,200]\) with the
bound \(\SstarLo\), and 154 cells, of width about \(0.007\) near the maximum,
cover \([\cStarLo,\cStarHi]\) with the bound \(\SstarHi\).  For \(c\ge200\),
\(\Gam(c)<2/3\) and \(\Apr(c)\ge\Apr(200)>221.77\) give
\(S(c)<F_{\Apr(200)}(2/3)<\SstarLo\).
\end{proof}

In Lean: \lean{optimum_global}, from \lean{savingF_pruned_lt_outside},
\lean{savingF_pruned_lt_global_tight}, \lean{mem_Ioo_of_savingF_pruned_ge} and
\lean{SstarPruned_bounds}; the supremum is \lean{SstarPruned}.  The cells are
generated by the script \texttt{search/certs.py} of the repository, which
chooses the endpoints and the points \(\theta\) in floating point and checks
every emitted inequality in exact rational arithmetic; Lean checks them again
(files \lean{Optimum/GlobalCells1.lean} to \lean{Optimum/GlobalCells5.lean}).  The
enclosures of the logarithms are those of \cite{upstream}, a series for
\(\ln\frac{1+z}{1-z}\) with an explicit remainder (\lean{log_mem}).

The same method certifies the optimum with the paper's encodings, which is the
optimum of Theorem~\ref{thm:T4}.

\begin{theorem}[The optimum over the ratio, paper's encodings]\label{thm:global-full}
Let \(S(c)=F_{\Afull(c)}(\Gam(c))\) be the supremum of the saving with the
paper's encodings at the ratio \(c\).  Then
\[
  \SfullLo<\sup_{c>10}S(c)\le\SfullHi,
\]
and every \(c>10\) with \(S(c)\ge\SfullLo\) lies in \((\cFullLo,\cFullHi)\).
\end{theorem}

\begin{proof}
As for Theorem~\ref{thm:global}, with \(\Afull\) in place of \(\Apr\): the lower
bound is Proposition~\ref{prop:values}(i), and 261 cells cover \((10,200]\),
87 below \(\cFullLo\) and 86 above \(\cFullHi\) with the bound \(\SfullLo\), and 88
on \([\cFullLo,\cFullHi]\) with the bound \(\SfullHi\); for \(c\ge200\),
\(\Afull(200)>238.74\).
\end{proof}

In Lean: \lean{SstarFull}, \lean{SstarFull_bounds}, \lean{SstarFull_le},
\lean{mem_Ioo_of_savingF_full_ge}, \lean{savingF_full_lt_global_tight},
\lean{lt_SstarFull_iff}, and the cells \lean{Optimum/GlobalFullCells1.lean} to
\lean{Optimum/GlobalFullCells6.lean}.

The upper bounds \(\SstarHi\) and \(\SfullHi\) are not strict, and the maximizers
are not located more precisely than the two windows: both functions are flat
near their maxima, \(S(21)\approx0.002093\) and \(S(23)\approx0.002092\) with
pruned encodings, and \(S(\cFullLo)\approx S(\cFullHi)\approx0.0020588\) with the
paper's encodings, just below \(\SfullLo\).

%% LEAN NAMES of the four theorems (file ImprovedExponents/Statements.lean):
%%   \lean{threeSum_paper}    Theorem T2 (the paper's exponents, better parameters)
%%   \lean{threeSum_exact}    Theorem T3 (+ exact count)
%%   \lean{threeSum_full}     Theorem T4 (+ true middle size)
%%   \lean{threeSum_pruned}   Theorem T1 (+ pruned encodings)
%% with exactTriangle_*, minPlus_*, apsp_* for the other problems, the savings
%% allSolved_*_num, and the general theorems allSolved_paper, allSolved_exact,
%% allSolved_full (Pipeline/General.lean), allSolved_pruned (Pipeline/PrunedClaim.lean).

\section{Consequences for 3SUM and APSP}
\label{sec:consequences}

The theorems of this section are statements about programs of the word RAM, in
the sense of Section~\ref{sec:machine}.  Four are for Exact Triangle and its
consequences; they add the ingredients one at a time, and the first adds none.
The fifth is the all-edges route to APSP.

\subsection{From Exact Triangle to the three problems}

\begin{proposition}[\cite{upstream}]\label{prop:transfer}
Let \(0\le\delta\le1\), and suppose that a program solves Exact Triangle on
\(n\) vertices per part with weights up to \(n^{\nu}\) in time
\(O(\nu\,n^{3-\delta}(\log n)^{e})\), uniformly in \(\nu\).  Then
\lean{ThreeSum.SolvedInTime}~\(a\) holds for every rational \(a>2-\delta/2\), and
\lean{MinPlusProduct.SolvedInTime}~\(a\) and \lean{APSP.SolvedInTime}~\(a\) hold for
every rational \(a>3-\delta/3\).
\end{proposition}

These are Theorem~21(a),(b) of the paper as formalized upstream
(\lean{threeSum_of_uniform}, \lean{minPlus_polylog_of_uniform},
\lean{apsp_polylog_of_uniform}).  The hypothesis is obtained from the bound of
Theorem~\ref{thm:total}, which holds from some size on, by solving smaller
instances by brute force (\lean{exactTriangleUniform_of_explicitFrom}).

\subsection{The regime test}
\label{sec:regime}

The upstream solver of the thin product tests whether \(D^{18}\le N\); it runs
the recursion if so and brute force if not.  The test is part of the program
text, and it restricts the ratios at which the program can be used.
\begin{itemize}
\item The recursion must be applicable in the whole regime \(D^{18}\le N\): the
  tile \(\sqrt K N_0\) must fit into \(N\).  For all large \(m\) this holds if
  \(\Apr(c)<18\ln 4\), that is, for \(c<\cFitsMax\ldots\)
  (\lean{exists_fits18}).
\item The instances with \(D\le N^{\varepsilon}\) must lie in the regime:
  \(\varepsilon\le1/18\).  If \(\Afull(c)\ge18\ln 4\), that is, for
  \(c\ge\cFullMin\ldots\), then \(\Rc(\gamma)\le\Rc(0)\le1/18\), and every
  admissible \(\varepsilon\) qualifies (\lean{thinR_le_inv}).
\end{itemize}
Both conditions hold at \(c=\cUncond\) (\lean{basePruned_108_5_lt},
\lean{baseFull_ge_18}), and the numerals of Theorems~\ref{thm:T2}
and~\ref{thm:T3}, which use the program as it is, are obtained there.  The
optimum \(c\approx\cFullApprox\) of the paper's encodings lies just below this
range: there \(\varepsilon\) has to exceed \(1/18\) slightly, while
\(\Apr(c)/\ln4\approx17.6\).

For Theorem~\ref{thm:T4} we therefore add a test with a rational exponent.
Three procedures are appended to the upstream program (\lean{programRS}): a
test \(D^{r}\le N^{s}\) for natural numbers \(r,s\) (\lean{testRSBody}: it
computes \(N^{s}\) by a loop and decides \(D^{r}<N^{s}+1\) with upstream's
routine \lean{powLtBody}, never forming a number above \(N^{s}D+N^{s}+D+r+s\)),
and a copy of the dispatching procedure that calls it.  The analysis of the
solver is carried out once for an abstract regime (\lean{RegimeTest},
\lean{thinClaimR_of_offlineWithin}), of which the upstream test and the new one
are two instances (\lean{regime18}, \lean{regimeRS}).  With the new test the two
conditions above become \(\Apr(c)<(r/s)\ln4\) and \(\varepsilon\le s/r\)
(\lean{exists_fitsRS}, \lean{thinClaimRS_gammaX}).  Such a fraction exists for
every admissible \(\varepsilon\): \(\varepsilon<\Rc(\gamma)\) means
\(1/\varepsilon>(\Afull(c)+\gamma\ln4)/\ln4>\Apr(c)/\ln4\), since
\(\Apr\le\Afull\) (\lean{basePruned_le_baseFull}).  The only remaining
condition on the ratio is \(\varepsilon\le1/16\), needed by the arithmetic of
Theorem~\ref{thm:total}; \(\Afull(c)\ge16\ln4\) for \(c\ge20\) gives it
(\lean{baseFull_ge_16}).

\subsection{The four theorems}

\begin{theorem}[The paper's exponents with better parameters]\label{thm:T2}
With the program, the reduction and the leaf count of the paper, at
\(c=\cUncond\) and suitable rational \(\theta\), \(\varepsilon\) and \(\eta\):
Exact Triangle is solved in time \(O(n^{3-\dTwo})\).  Consequently
\lean{ThreeSum.SolvedInTime} holds with the exponent \(\ThreeSumTwo\), and
\lean{MinPlusProduct.SolvedInTime} and \lean{APSP.SolvedInTime} with the
exponent \(\APSPTwo\).
\end{theorem}

Nothing is new here except the choice of parameters.  Corollary~26 of the paper
rounds \(\varepsilon\) down to \(1/18\) and takes \(c=21\); with
\(\gamma=\min(\gpap(c,\theta),\frac12-q(\theta))\) balanced and
\(\varepsilon\) close to \(\Rc(\gamma)\) the saving
\(\varepsilon\gamma/2\) exceeds \(\dTwo\) (\lean{t2_witness}).  In Lean: the general
statement is \lean{allSolved_paper}, from \lean{thinClaim_gammaOf},
\lean{hostBound_rat} and \lean{explicitFrom_original}; the saving is
\lean{allSolved_paper_num}, and the end statements are
\lean{exactTriangle_paper}, \lean{threeSum_paper}, \lean{minPlus_paper} and
\lean{apsp_paper}.

\begin{theorem}[With the exact count]\label{thm:T3}
With the same program, analyzed by Theorem~\ref{thm:offline}: Exact Triangle is
solved in time \(O(n^{3-\dThree})\).  Consequently
\lean{ThreeSum.SolvedInTime} holds with the exponent \(\ThreeSumThree\), and
\lean{MinPlusProduct.SolvedInTime} and \lean{APSP.SolvedInTime} with the
exponent \(\APSPThree\).
\end{theorem}

In Lean: \lean{allSolved_exact}, from \lean{thinClaim_gammaX} and the same
reduction; the witness is \lean{t3_witness}, the saving
\lean{allSolved_exact_num}, and the end statements are
\lean{exactTriangle_exact}, \lean{threeSum_exact}, \lean{minPlus_exact} and
\lean{apsp_exact}.

\begin{theorem}[With the true middle size]\label{thm:T4}
Let \(c\ge20\) and \(0\le\delta<F_{\Afull(c)}(\Gam(c))\).  With the host of
Theorem~\ref{thm:mid}, the thin product of Theorem~\ref{thm:offline} and the
regime test with a rational exponent, Exact Triangle is solved in time
\(O(n^{3-\delta})\); consequently so is every \(\delta\) below
\(\sup_{c>10}F_{\Afull(c)}(\Gam(c))\in(\SfullLo,\SfullHi]\).  At
\(c=\cUncondOpt\) this holds for \(\delta=\SfullLo\).  Consequently
\lean{ThreeSum.SolvedInTime} holds with the exponent \(\ThreeSumFour\), and
\lean{MinPlusProduct.SolvedInTime} and \lean{APSP.SolvedInTime} with the
exponent \(\APSPFour\).
\end{theorem}

\begin{proof}
By Theorem~\ref{thm:sup} and Proposition~\ref{prop:rational} there is an
admissible rational tuple \((c',\theta',\sigma',\varepsilon')\) near \(c\) with
saving above \(\delta\).  Lower \(\gamma\) below \(\gX(c',\theta')\) and
\(\varepsilon\) to a rational so that the saving stays above \(\delta\)
(\lean{exists_gamma_rat_of_lt}), and choose \(r/s\) between
\(\Apr(c')/\ln4\) and \(1/\varepsilon'\) (Section~\ref{sec:regime}).
Theorem~\ref{thm:offline} with the new test gives the thin product, and
Theorem~\ref{thm:total} with \(\eta=1-1/(2\sigma)\) and
\(d=\varepsilon/(1-\eta)\) slightly lowered gives Exact Triangle.  Its second
condition holds because \(\frac32d\le3\sigma\varepsilon\) and
\(\delta<\varepsilon(1-\sigma)\), so
\(\frac32d+\delta<(1+2\sigma)\varepsilon<3\varepsilon\le\frac3{16}\).  The
supremum over \(c\) is Theorem~\ref{thm:global-full}, and the value at
\(c=\cUncondOpt\) is Proposition~\ref{prop:values}(i).
\end{proof}

In Lean: \lean{allSolved_full} is the statement for every \(c\ge20\),
\lean{allSolved_full_sup} the statement for every saving below the supremum
\lean{SstarFull}, \lean{allSolved_full_num} the saving at \(c=\cUncondOpt\),
and the end statements are \lean{exactTriangle_full}, \lean{threeSum_full},
\lean{minPlus_full} and \lean{apsp_full}.  The five-digit exponents are those
of \(\delta=\dFour\): the optimum improves the saving at \(c=\cUncond\) by less
than \(2\cdot10^{-7}\).

\begin{theorem}[With pruned encodings]\label{thm:T1}
Let \(c\ge21\) and \(0\le\delta<F_{\Apr(c)}(\Gam(c))\).  With the host of
Theorem~\ref{thm:mid} and the thin product of Theorem~\ref{thm:encoder}, Exact
Triangle is solved in time \(O(n^{3-\delta})\); consequently so is every
\(\delta\) below \(\sup_{c>10}F_{\Apr(c)}(\Gam(c))\in(\SstarLo,\SstarHi]\).  At
\(c=\cCond\) this holds for \(\delta=\dOne\).  Consequently
\lean{ThreeSum.SolvedInTime} holds with the exponent \(\ThreeSumOne\), and
\lean{MinPlusProduct.SolvedInTime} and \lean{APSP.SolvedInTime} with the
exponent \(\APSPOne\).
\end{theorem}

The proof is that of Theorem~\ref{thm:T4} with \(\RcP\) in place of \(\Rc\) and
Theorem~\ref{thm:encoder} in place of Theorem~\ref{thm:offline}; the test with
a rational exponent is part of the program \lean{programP}.  For \(c\ge21\) we
have \(\Apr(c)\ge16\ln4\) (\lean{basePruned_ge_16}), so \(\varepsilon\le1/16\)
and \(\frac32d+\delta<3\varepsilon\le0.19\).  The supremum over \(c\) is
Theorem~\ref{thm:global}.  At \(c=\cCond\) the parameters are
\(\Gam>\GamCond\), \(\theta\approx\thetaCond\), \(\sigma\approx\sigmaCond\) and
\(\varepsilon\) just below \(R>\epsCond\) (\lean{Gam_22_bounds},
\lean{thinR_pruned_bounds}, \lean{savingF_pruned_gt}).  In Lean:
\lean{allSolved_pruned} is the statement for every \(c\ge21\) (file
\lean{Pipeline/PrunedClaim.lean}; the form that still takes the thin-product
claim as an argument is \lean{allSolved_pruned_of_claim}, from
\lean{explicitFrom_pruned_of_claim}), \lean{allSolved_pruned_sup} the
statement for every saving below the supremum \lean{SstarPruned},
\lean{allSolved_pruned_num} the saving at \(c=\cCond\), and the end statements
are \lean{exactTriangle_pruned}, \lean{threeSum_pruned}, \lean{minPlus_pruned}
and \lean{apsp_pruned}.  No hypothesis is involved.

\subsection{APSP through all-edges Exact Triangle}
\label{sec:alledges}

Footnote~10 of the paper observes that APSP can keep all of the saving of Exact
Triangle instead of a third: "the proof of Theorem~17 also solves the
all-edges version of Exact Triangle (scanning each pair only until its zero
triangle is found), to which the \(\minplus\)-product reduces on the same \(n\)
vertices".  We carry this out as programs.  The \emph{all-edges version} of
Exact Triangle asks, for every pair \((a,b)\) of the two outer parts, whether
some \(c\) makes \(w(a,b)+w(b,c)+w(a,c)=0\); as a task of the machine
(\lean{aeTask}) it writes the \(n^{2}\) answers to an output array
(\lean{aeFlags}).

\begin{lemma}[Comparisons as equalities]\label{lem:F}
Let \(0\le x,y,u<2^{b}\) be integers and
\(f(\ell)=\lfloor u/2^{\ell}\rfloor-\lfloor x/2^{\ell}\rfloor-\lfloor y/2^{\ell}\rfloor\).
Then
\[
  x+y<u\iff f(0)=1\ \text{ or }\ f(\ell)\in\{2,3\}\text{ for some }\ell\le b .
\]
\end{lemma}

\begin{proof}
\(f(b)=0\), and \(f(\ell)=2f(\ell+1)+e_\ell\), where \(e_\ell\in\{-2,-1,0,1\}\)
is the bit of \(u\) at position \(\ell\) minus those of \(x\) and \(y\).  So
\(f(\ell+1)\ge2\) implies \(f(\ell)\ge2\).  If \(f(\ell)\in\{2,3\}\) for some
\(\ell\), then \(f(0)\ge2\) by descent, and \(f(0)=u-x-y\).  Conversely let
\(f(0)\ge1\).  If \(f(0)\ne1\), let \(\ell\) be the largest index with
\(f(\ell)\ge2\); then \(\ell<b\), \(f(\ell+1)\le1\) and
\(f(\ell)\le2f(\ell+1)+1\le3\).
\end{proof}

In Lean: \lean{add_lt_iff_shiftDiff}.

\begin{proposition}\label{prop:minplus-math}
Let \(A,B\) be \(n\times n\) integer matrices with entries of absolute value at
most \(W\), let \(2W<2^{b}\), and let \(t\) be a matrix of thresholds with
\(-2W\le t_{ij}<2^{b}-2W\).  For \(\ell\le b\) put
\(X_\ell(i,k)=\lfloor(A_{ik}+W)/2^{\ell}\rfloor\),
\(Y_\ell(k,j)=\lfloor(B_{kj}+W)/2^{\ell}\rfloor\) and
\(U_\ell(i,j)=\lfloor(t_{ij}+2W)/2^{\ell}\rfloor\).  Then, for every pair
\((i,j)\), there is a \(k\) with \(A_{ik}+B_{kj}<t_{ij}\) if and only if, for
one of the \(2b+3\) pairs \((\ell,r)\) with \((\ell,r)=(0,1)\) or
\(r\in\{2,3\}\), there is a \(k\) with
\[
  X_\ell(i,k)+Y_\ell(k,j)+\bigl(r-U_\ell(i,j)\bigr)=0 .
\]
Consequently a binary search that runs for all pairs \((i,j)\) in parallel
computes the \(\minplus\)-product of \(A\) and \(B\) from the answers to
\(R(2R+3)\) all-edges Exact Triangle instances on the same \(n\) vertices per
part, where \(R=\lceil\log_2(4W+1)\rceil\).
\end{proposition}

\begin{proof}
Apply Lemma~\ref{lem:F} to \(x=A_{ik}+W\), \(y=B_{kj}+W\) and
\(u=t_{ij}+2W\).  The entry \(C_{ij}\) of the product is the unique integer
\(c\in[-2W,2W]\) such that the answer is negative for the threshold \(c\) and
positive for \(c+1\); a binary search finds it in \(R\) rounds, and a round
asks one matrix of thresholds.
\end{proof}

In Lean: \lean{exists_lt_iff_exists_exactTriangle}, \lean{minPlus_unique},
\lean{search_eq_minPlus} and \lean{search_exactTriangle}.

The reduction of the paper (Theorem~21(b), \cite{VW18}) ends with a binary search of
exactly this kind, over a procedure that answers the threshold questions for
all pairs at once (\lean{pairsTask}, \lean{isHost_mp}), followed by
repeated squaring for APSP (\lean{claim_apspFromMinPlus}).  The factor \(3\) is
lost before that, in a blocking of size \(n^{1/3}\) that reduces the all-pairs
questions to the decision problem.  Two hosts replace the blocking.

\begin{theorem}[All pairs from all-edges Exact Triangle]\label{thm:pairs}
There is a program that, given a procedure solving all-edges Exact Triangle
within \(T(n,U)\) steps, answers the threshold questions of
Proposition~\ref{prop:minplus-math} for all pairs within
\((2b+3)\bigl(T(n,9U)+O(n^{2})\bigr)+O(n^{2})\) steps, where
\(b=\lfloor\log_2 3U\rfloor+1\); in particular within
\(O\bigl((1+\log U)(T(n,9U)+n^{2})\bigr)\).
\end{theorem}

\begin{proof}
With \(W=U\) the hypotheses of Proposition~\ref{prop:minplus-math} hold, and
the weights of the instances are at most \(2^{b}+3\le9U\).  The program writes
\(A+U\), the transpose of \(B+U\) and \(t+2U\), and then runs the \(b\) levels
from the top: at each level it asks the two instances with \(r=2,3\) and takes
the disjunction of their answers into the output, then halves every number
(upstream's pass for \(\lfloor x/2^{\ell}\rfloor\), which avoids division); at
level \(0\) it also asks \(r=1\).  The output is then the matrix of answers.
\end{proof}

In Lean: \lean{isHost_pairsAE} with the time \lean{pairsTimeAE}, from the
bridge \lean{pairFlags_iff} between the two tasks and the invariant
\lean{orFlags_eq_pairFlags}; the bound \lean{pairsTimeAE_le} with the constant
\lean{CA}, and the transfer \lean{pairsFromAllEdges}.

\begin{theorem}[Theorem~\ref{thm:mid} for the all-edges version]\label{thm:mid-ae}
Theorem~\ref{thm:mid} holds for all-edges Exact Triangle in place of Exact
Triangle, with the same bound and another constant \(C\).
\end{theorem}

\begin{proof}
The host of Theorem~17 keeps \(n^{2}\) flags, one per pair \((a,b)\), zero at
the start.  For every query pair that an instance accepts and whose flag is
still \(0\), it scans the piece \(C_k\) of the instance and stores the result
in the flag; pairs whose flag is \(1\) are skipped.  Every pair with a zero
triangle is accepted, and hit, in the instance of the piece of its \(c\), so
the flags are the answers at the end.  A scan either fails, which the paper
charges to a false positive of the prime, or turns a flag from \(0\) to \(1\);
so there are at most \(F(p)+n^{2}\) scans.  The \(n^{2}\) additional scans cost
\(O(n^{2}\sqrt{D'}/g')\), which is within the term \(n^{2}D'g'\); zeroing and
copying the flags costs \(O(n^{2})\).
\end{proof}

In Lean the pure part is \lean{flagsAt_m} (the flags after all instances are
the answers) and \lean{sum_execsAE_le} (the count of the scans); the
procedure that scans the accepted pairs is \lean{scanPairsAEBody}, the brute
force of the small case \lean{bruteAEBody}, the host \lean{ae17_solves}, its
time \lean{hostTimeAE} and \lean{obeysBound17Mid_hostTimeAE}, and the claim
\lean{claim17MidAE_of_pack}.  The statement of Theorem~\ref{thm:mid} is the
same proposition \lean{Claim17Mid} in a time model (\lean{lightModelAE}) whose
reading of "Exact Triangle" is the all-edges task; the deductions of
Sections~\ref{sec:middle} and~\ref{sec:consequences} are made for an arbitrary
model, so they apply unchanged (\lean{MidHostRat}, \lean{midHostRat_ae},
\lean{explicitFrom_full}, \lean{explicitFrom_pruned}).

\begin{theorem}[APSP through all-edges Exact Triangle]\label{thm:apsp-ae}
Let \(c\ge20\) and \(0\le\delta<F_{\Afull(c)}(\Gam(c))\), or let \(\delta\) be
any number below \(\sup_{c>10}F_{\Afull(c)}(\Gam(c))>\SfullLo\).  Then the
\(\minplus\)-product and APSP are solved in time \(O(n^{a})\) for every
rational \(a>3-\delta\); consequently \lean{MinPlusProduct.SolvedInTime} and
\lean{APSP.SolvedInTime} hold with the exponent \(\APSPFn\).  With pruned
encodings the same holds for \(c\ge21\) and \(\delta<F_{\Apr(c)}(\Gam(c))\),
for every \(\delta<\sup_c F_{\Apr(c)}(\Gam(c))>\SstarLo\), and with the
exponent \(\APSPFnOne\).
\end{theorem}

\begin{proof}
Theorems~\ref{thm:T4} and~\ref{thm:T1} with Theorem~\ref{thm:mid-ae} in place
of Theorem~\ref{thm:mid} give all-edges Exact Triangle in time
\(O(n^{3-\delta}\log n)\) for weights up to \(n^{O(1)}\).  Theorem~\ref{thm:pairs}
and the search of Proposition~\ref{prop:minplus-math} (upstream's program,
\(O(\log U)\) rounds at the bound \(6U\)) give the \(\minplus\)-product of
matrices with entries up to \(u\) in time \(O(\log^{2}u\,(T(n,54u)+n^{2}))\),
and repeated squaring gives APSP with \(O(\log n)\) products of matrices with
entries up to \(O(n\,u)\).  Along \(u=n^{\kappa}\) every factor is
\(n^{3-\delta}\) times a polylogarithm.
\end{proof}

In Lean: \lean{minPlusSolved_full}, \lean{minPlusSolved_full_sup},
\lean{minPlusSolved_pruned} and \lean{minPlusSolved_pruned_sup} (the structure
\lean{MinPlusSolved}), through \lean{solvedIn_mp_of_allEdges},
\lean{minPlusInPolylog_of_explicitFrom}, \lean{apspInPolylog_of_explicitFrom}
and \lean{minPlusSolved_of_explicitFrom}; the end statements are
\lean{minPlus_allEdges}, \lean{apsp_allEdges}, and with pruned encodings
\lean{minPlus_allEdges_pruned}, \lean{apsp_allEdges_pruned}, the strongest
statements of the paper for the \(\minplus\)-product and APSP.

\section{The formalization}
\label{sec:formalization}

The Lean development is the directory \lean{ImprovedExponents} of the
repository \url{https://github.com/qawbecrdtey/ImprovedExponents}, which also
contains this paper.  It has \LeanFiles{} files and
\LeanLines{} lines, of which \LeanCellLines{} are generated certificates for
Theorems~\ref{thm:global}, \ref{thm:global-full} and~\ref{prop:shapes}; about
\LeanHandLines{} lines are written by hand.

\subsection{Relation to the upstream formalization}

The part of the formalization \cite{upstream} of the paper that we use, the
\UpstreamFiles{} files of Lean that our development imports directly or
transitively, is included in the repository at a fixed revision
(\texttt{\UpstreamRev}; directory \texttt{upstream/3sum-apsp}) and ported
from Lean and Mathlib \texttt{\UpstreamVersion}, for which it was written, to
\texttt{\LeanVersion}.  The port changes only what no longer compiles, in
\UpstreamChanged{} of the files, and no statement or definition; every other
file is upstream's, byte for byte.  It provides the word RAM
and the notion \lean{Problem.SolvedInTime} (file \lean{EndStatement.lean}), a
small programming language with a time model (\lean{lightModel}) and a compiler
to the word RAM, the programs of the paper, and the proofs of its theorems.

We reuse the following.
\begin{itemize}
\item The program of Section~4, \lean{program31}, for arbitrary rational
  parameters, with its correctness and the bound of its running time in terms
  of the two costs of Theorem~30.
\item The count of the boxes, \lean{lemma_29_count}, which leaves the tail
  \(\sum_{d\ge t}\beta_d\) explicit.
\item The host of Theorem~17 and the passage from a host to a claim about
  programs, \lean{claim17_of_host}.
\item The reductions of Theorem~21 and the passage to the end statements:
  \lean{threeSum_of_uniform}, \lean{minPlus_polylog_of_uniform},
  \lean{apsp_polylog_of_uniform} and \lean{SolvedInTime.endStatement}.
\item The enclosures of logarithms between rational numbers, \lean{log_mem}.
\end{itemize}
Upstream proves the bounds of the paper at the parameters of the paper.  Where
an upstream proof fixes parameters, we repeat it for general ones: the solver
of the thin product at arbitrary ratios (\lean{Pipeline/ThinClaim.lean}), the
time of the program against the exact tail (\lean{Cost8X}), the routines that
compute the parameters of the host (\lean{ParamRoutines}), and the arithmetic
of the total time (\lean{Total}).  Where an upstream proof assumes the full
encodings, we repeat it under the weaker contract of
Section~\ref{sec:encoder-program}: the encoder of Theorem~5, the procedures of
Theorem~30 and the routines with chosen parameters (\lean{PrunedProgram}), and
the time of the new program against \lean{cost8P} (\lean{Cost8P}).  The file
\texttt{NOTICE} of the repository lists the files that are adapted from
upstream, which is distributed under the Apache~2.0 license.

\subsection{What is new}

Table~\ref{tab:dirs} lists the directories.

\begin{table}[ht]
\centering\small
\begin{tabular}{lr>{\raggedright\arraybackslash}p{3.3in}}
\toprule
directory & lines & content\\
\midrule
\lean{ExactCount} & \LeanExactCountLines & \(\gtwo\), Theorem~\ref{thm:tail},
  Proposition~\ref{prop:lemmaA}\\
\lean{Cost8X} & \LeanCostEightXLines & the upstream program against the exact tail
  (Theorem~\ref{thm:offline})\\
\lean{PrunedEncoding} & \LeanPrunedEncodingLines & Theorem~\ref{thm:pruned},
  Proposition~\ref{prop:costP}\\
\lean{PrunedProgram} & \LeanPrunedProgramLines & the pruned encoder and the
  solver under the weaker contract (Section~\ref{sec:encoder-program})\\
\lean{Cost8P} & \LeanCostEightPLines & the new program against \lean{cost8P}
  (Theorem~\ref{thm:encoder})\\
\lean{HostMid} & \LeanHostMidLines & the host with one added assignment
  (Theorem~\ref{thm:mid})\\
\lean{ParamRoutines} & \LeanParamRoutinesLines & routines for \(\lfloor n^{a/b}\rfloor\) and
  \(\lceil D^{c/d}\rceil\)\\
\lean{Pipeline} & \LeanPipelineLines & thin-product claims, the regime
  tests, the offline routines, the general theorems\\
\lean{Total} & \LeanTotalLines & Theorem~\ref{thm:total}\\
\lean{Optimum} & \LeanOptimumLines & Section~\ref{sec:optimum},
  Proposition~\ref{prop:half}, the shapes of Section~\ref{sec:limits};
  \LeanCellLines{} lines are generated\\
\lean{MinPlus} & \LeanMinPlusLines & Lemma~\ref{lem:F},
  Proposition~\ref{prop:minplus-math}\\
\lean{AllEdges} & \LeanAllEdgesLines & the all-edges route to APSP
  (Theorems~\ref{thm:pairs}, \ref{thm:mid-ae}, \ref{thm:apsp-ae})\\
\lean{Statements.lean} & \LeanStatementsLines & the theorems of
  Section~\ref{sec:consequences}\\
\bottomrule
\end{tabular}
\caption{The Lean development.}
\label{tab:dirs}
\end{table}

Two points of the formalization differ from a proof on paper.

\emph{The programs are almost fixed.}  In Theorems~\ref{thm:T2}
and~\ref{thm:T3} the program for the thin product and the host of the
reduction are upstream's, at other parameters; the only new program text is
that of the routines that compute \(D'=\lfloor n^{a/b}\rfloor\) and
\(g'=\lceil D'^{c/d}\rceil\).  So the improvement from \(\ETpaper\) to
\(\dThree\) is obtained by a new analysis of existing programs.
Theorem~\ref{thm:T4} adds one assignment to the host and three short
procedures for the regime test; the all-edges route (Theorem~\ref{thm:apsp-ae})
adds the two hosts of Section~\ref{sec:alledges}, which reuse upstream's
passes, scans, hashing and instance construction.  Theorem~\ref{thm:T1} adds
seven procedures (\lean{programP}): the pruned encoder and six that call it,
or call upstream's procedures under the new contract; the procedures that
build the tiles and answer the queries are upstream's, with a new proof.

\emph{Numerals are corollaries.}  The general theorems
(\lean{allSolved_paper}, \lean{allSolved_exact}, \lean{allSolved_full},
\lean{allSolved_pruned}) quantify over the parameters or over all savings
below the closed form \(S(c)\).  The decimal exponents of
Table~\ref{tab:results} follow from certified enclosures of \(S(c)\) at one
ratio (Proposition~\ref{prop:values}).

\subsection{The main statements}

Table~\ref{tab:names} lists the Lean names of the main statements.

\begin{table}[ht]
\centering\footnotesize
\setlength{\tabcolsep}{4pt}
\begin{tabular}{lll}
\toprule
statement & declaration & file\\
\midrule
Theorem~\ref{thm:tail} & \lean{sum_beta_le_exp} & \lean{ExactCount/Tail.lean}\\
Proposition~\ref{prop:lemmaA} & \lean{sum_card_Leaves_le_exp}
  & \lean{ExactCount/Leaves.lean}\\
Theorem~\ref{thm:offline} & \lean{offlineWithinX} & \lean{Cost8X/Offline.lean}\\
Theorem~\ref{thm:pruned} & \lean{pruned_encoding}
  & \lean{PrunedEncoding/Pruned.lean}\\
Proposition~\ref{prop:costP} & \lean{costsP} & \lean{PrunedEncoding/Costs.lean}\\
Remark~\ref{rem:encwork} & \lean{encWorkP_le_M_div} & \lean{PrunedProgram/Work.lean}\\
Theorem~\ref{thm:encoder} & \lean{prunedEncoderClaim}
  & \lean{Pipeline/PrunedClaim.lean}\\
 & \lean{encodeP_spec}, \lean{tEncP_le}
  & \lean{PrunedProgram/Encode/Recursion.lean}\\
 & \lean{offline32P_meets} & \lean{PrunedProgram/Wrapper/Offline.lean}\\
 & \lean{offlineWithinP} & \lean{Cost8P/Offline.lean}\\
Theorem~\ref{thm:mid} & \lean{claim17Mid_of_paramProcs}
  & \lean{HostMid/AtParameters.lean}\\
Theorem~\ref{thm:total} & \lean{explicitFrom_mid} & \lean{Total/Explicit.lean}\\
Theorem~\ref{thm:sup} & \lean{isLUB_tupleSaving}
  & \lean{Optimum/PerCSupremum.lean}\\
Proposition~\ref{prop:rational} & \lean{exists_rat_tuple_basePruned}
  & \lean{Optimum/PerCDensity.lean}\\
Proposition~\ref{prop:half} & \lean{isLUB_paperSaving},
  \lean{savingF_GamHalf_lt} & \lean{Optimum/PaperAccounting.lean}\\
Theorem~\ref{thm:global} & \lean{optimum_global} & \lean{Statements.lean}\\
Theorem~\ref{thm:global-full} & \lean{optimum_global_full} & \lean{Statements.lean}\\
Theorem~\ref{thm:T2} & \lean{threeSum_paper}, \lean{apsp_paper}
  & \lean{Statements.lean}\\
Theorem~\ref{thm:T3} & \lean{threeSum_exact}, \lean{apsp_exact}
  & \lean{Statements.lean}\\
Theorem~\ref{thm:T4} & \lean{allSolved_full}, \lean{threeSum_full}
  & \lean{Statements.lean}\\
Theorem~\ref{thm:T1} & \lean{allSolved_pruned} & \lean{Pipeline/PrunedClaim.lean}\\
 & \lean{threeSum_pruned}, \lean{apsp_pruned} & \lean{Statements.lean}\\
Lemma~\ref{lem:F} & \lean{add_lt_iff_shiftDiff}
  & \lean{MinPlus/Comparison.lean}\\
Proposition~\ref{prop:minplus-math} & \lean{search_exactTriangle}
  & \lean{MinPlus/ExactTriangle.lean}\\
Theorem~\ref{thm:pairs} & \lean{isHost_pairsAE} & \lean{AllEdges/Pairs/Host.lean}\\
Theorem~\ref{thm:mid-ae} & \lean{ae17_solves} & \lean{AllEdges/Host/ProgramAE.lean}\\
 & \lean{midHostRat_ae} & \lean{AllEdges/Host/ClaimAE.lean}\\
Theorem~\ref{thm:apsp-ae} & \lean{minPlusSolved_full}, \lean{apsp_allEdges}
  & \lean{Statements.lean}\\
 & \lean{minPlusSolved_pruned}, \lean{apsp_allEdges_pruned}
  & \lean{Statements.lean}\\
Proposition~\ref{prop:shapes} & \lean{shapes_table} & \lean{Optimum/Shapes/Table.lean}\\
\bottomrule
\end{tabular}
\caption{The main statements and their Lean names.  The end statements for
Exact Triangle, the \(\minplus\)-product and APSP are named in the same way
(\texttt{exactTriangle\_full}, \texttt{minPlus\_full}, \texttt{apsp\_full},
\texttt{minPlus\_allEdges}, and so on).}
\label{tab:names}
\end{table}

\subsection{What has to be trusted}

For the end statements of Theorems~\ref{thm:T2}, \ref{thm:T3}, \ref{thm:T4},
\ref{thm:T1} and~\ref{thm:apsp-ae}: the Lean kernel, and the definitions of
upstream's \lean{EndStatement.lean}, which state what it means that a problem is
solved in time \(O(n^{r})\) on the word RAM.  That file imports nothing, and
the module \texttt{ImprovedChallenge.lean} of the repository, which states the
four strongest end statements, contains its 132 lines of Lean verbatim and
imports nothing either.  An end statement such as
\begin{center}
\leanraw{theorem threeSum_pruned : EndStatement.ThreeSum.SolvedInTime 1.99896}
\end{center}
mentions no definition of ours and takes no argument.  No statement of this
paper rests on a hypothesis.

The development contains no \texttt{sorry} and declares no axiom.  Lean
reports that the end statements depend on the axioms \texttt{propext},
\texttt{Classical.choice} and \texttt{Quot.sound}, the three axioms that
Mathlib uses throughout, and on no other; in particular the numerical
certificates are checked by the kernel and not by compiled code.

\subsection{How to check}

In the root of the repository, \texttt{lake exe cache get} fetches the compiled
Mathlib, and \texttt{lake build} checks everything: the included part of
upstream, the library and the two modules described next; the script
\texttt{scripts/check-upstream.sh} compares the included part with upstream's
revision and confirms that every changed file is marked and listed.  The module \texttt{ImprovedChallenge.lean}
states the four strongest end statements, \lean{threeSum_pruned}
(\(\ThreeSumOne\)), \lean{exactTriangle_pruned} (\(\ETOne\)),
\lean{minPlus_allEdges_pruned} and \lean{apsp_allEdges_pruned}
(\(\APSPFnOne\); Theorems~\ref{thm:T1} and~\ref{thm:apsp-ae}), with
upstream's \lean{EndStatement.lean} copied into it verbatim (the script
\texttt{scripts/check-endstatement.sh} confirms the copy against the included
file, which is upstream's) and nothing imported, each with \texttt{sorry} for
its proof; the module \texttt{ImprovedSolution.lean} imports the proofs.  The
two modules and the configuration \texttt{comparator.json} are in the form
that the tool Comparator \cite{comparator} checks: that the solution proves
the statements of the challenge, literally, that it uses only the axioms
\texttt{propext}, \texttt{Quot.sound} and \texttt{Classical.choice}, and that
everything replays in Lean's kernel and in independent kernels; the
compilations run in a sandbox.  From Lean \texttt{v4.35.0-rc2} on, Comparator
ships with the toolchain as \texttt{lake comparator}, with the independent
kernels NanoDa and con-ron.  The script \texttt{scripts/verify-comparator.sh}
runs it as the Palomar registry does, in a bubblewrap sandbox and in a fresh
copy of the sources without build products, so that the challenge, the
solution and everything they import are compiled inside the sandbox.  We ran it
on the final development: \texttt{lake comparator} accepted the pair on
\ComparatorDate{} (``Your solution is okay!'', all three kernels; the log is
\texttt{scripts/comparator-run.log}).  Every Lean
file of the repository uses Lean's module system, and the statement module
imports nothing, as the Palomar registry \cite{Palomar} requires of a
submission.

The certificates of Theorems~\ref{thm:global} and~\ref{thm:global-full} are
regenerated by \texttt{python3 -m search.certs}, those of
Proposition~\ref{prop:shapes} with the option \texttt{--encoding shapes}, and
the script
\texttt{check\_lean\_names.py} next to the source of this paper checks that
every Lean name cited here exists in the sources.

\subsection{What is not formalized}

Only parts of Section~\ref{sec:limits}: the discussion of mixtures of
shapes, the grid of hypothetical shapes, the experiments on the arrangement of
the wanted set, and the remark on \(\omega\).  Each is marked ``\nf'' there.
Everything else, including the comparison of the shapes of Sch\"onhage's
family (Proposition~\ref{prop:shapes}), is machine-checked.

\section{Limits of the method}
\label{sec:limits}

The saving \(\SstarLo\) is the end of what the parameters give.  This section
records why the base identity cannot be exchanged for a better one of the same
kind, and what would have to change.  The comparison of the shapes of
Sch\"onhage's family is certified (Proposition~\ref{prop:shapes}): it is a
statement about the saving formula of Sections~\ref{sec:exact}
to~\ref{sec:optimum} evaluated for other shapes, since the programs exist for
the shape of the paper only.  The rest of the section, the mixtures, the
hypothetical shapes, the experiments and the remark on \(\omega\), is \nf; its
numbers are computed in floating point by the scripts in the directory
\texttt{search} of the repository.

\subsection{The family of Sch\"onhage}

For \(k,n\ge2\), Sch\"onhage's identity for
\(\schon k1n\oplus\schon1{(k-1)(n-1)}1\) has \(kn+1\) terms \cite{Sch81}.  It
has \(s=kn\) outer outputs with one term each, and one shared term for an
inner product of length \(b=(k-1)(n-1)\).  The paper uses \(k=n=3\): \(s=9\),
\(b=4\), ten terms.  The analysis of Sections~\ref{sec:exact}
to~\ref{sec:optimum} is written for this shape.  For another member of the
family it would be repeated with the corresponding counts of terms, variables
and leaves, an asymmetric shape being used together with its transpose so that
tiles stay square; we do not write the programs for other shapes, and what
follows is the saving formula of Theorem~\ref{thm:sup} evaluated for them.
With \(s\) outer outputs and inner length \(b\) one has \(D=b^{m}\),
\(N_0=\sqrt s^{\,L-m}\),
\(\gtwo\) with \(\delta\ln s\) in place of \(\delta\ln 9\), \(\gX\) and \(q\)
with \(\ln b\) in place of \(\ln 4\), the window \((0,s/(s+1))\) for \(\theta\)
in place of \((0,0.9)\), and the thinness constant
\(\Apr(c)=\frac c2\Hent(\frac1c)+(c-1)\frac{\ln s}2\); the saving at the ratio
\(c>s+1\) is then \(S(c)=F_{\Apr(c)}(\Gam(c))\) as in Theorem~\ref{thm:sup}.
In Lean (directory \lean{Optimum/Shapes}): \lean{g2S}, \lean{gammaXS},
\lean{qS}, \lean{gammaQS}, \lean{GamS}, \lean{thinRS}, \lean{savingS},
\lean{basePrunedS} and the saving \lean{shapeSaving}~\(s\,b\,c\); for
\((s,b)=(9,4)\) they reduce to the functions of Sections~\ref{sec:exact}
to~\ref{sec:optimum} (\lean{shapeSaving_nine_four}, from \lean{g2S_nine},
\lean{gammaXS_nine_four}, \lean{qS_nine_four}, \lean{GamS_nine_four} and
\lean{basePrunedS_nine}).  Lemma~\ref{lem:crossing} and the monotonicity of
\(\Apr\) carry over under \(2\le b<(s+1)^{2}\) and \(c>s+1\)
(\lean{GamS_pos}, \lean{GamS_lt_two_thirds}, \lean{GamS_monotoneOn},
\lean{basePrunedS_monotoneOn}).  Table~\ref{tab:family} gives the optimum for
each shape, with pruned encodings and the prime at its fixed point.

\begin{table}[ht]
\centering\small
\setlength{\tabcolsep}{4.5pt}
\begin{tabular}{lrrlrrr}
\toprule
shape \((k,n)\) & \(s\) & \(b\) & \(\sup_{c}S(c)\) (certified)
  & best \(c\) & \(\Gam\) & \(\varepsilon\)\\
\midrule
\((3,3)\) & 9 & 4 & \((\SstarLo,\SstarHi]\) & 21.92 & 0.0759 & 0.0552\\
\((2,4)\) & 8 & 3 & \((0.001875,0.001877]\) & 20.11 & 0.0749 & 0.0501\\
\((3,4)\) & 12 & 6 & \((0.001867,0.001869]\) & 27.96 & 0.0746 & 0.0501\\
\((2,5)\) & 10 & 4 & \((0.001776,0.001778]\) & 24.24 & 0.0742 & 0.0479\\
\((2,3)\) & 6 & 2 & \((0.001669,0.001671]\) & 16.03 & 0.0742 & 0.0450\\
\((2,6)\) & 12 & 5 & \((0.001623,0.001625]\) & 28.39 & 0.0733 & 0.0443\\
\((3,5)\) & 15 & 8 & \((0.001613,0.001615]\) & 34.04 & 0.0732 & 0.0441\\
\((4,4)\) & 16 & 9 & \((0.001573,0.001575]\) & 35.96 & 0.0729 & 0.0432\\
\((4,5)\) & 20 & 12 & \((0.001323,0.001325]\) & 44.01 & 0.0714 & 0.0371\\
\((5,5)\) & 25 & 16 & \((0.001098,0.0011]\) & 54.02 & 0.0698 & 0.0315\\
\bottomrule
\end{tabular}
\caption{The saving for Exact Triangle with Sch\"onhage's identity of shape
\((k,n)\) and pruned encodings.  The enclosures of the suprema over the ratio
\(c>s+1\) are certified in Lean (Proposition~\ref{prop:shapes};
Theorem~\ref{thm:global} for \((3,3)\)).  The last three columns, the ratio at
which the supremum is approached and the parameters there, are computed in
floating point and are not certified.}
\label{tab:family}
\end{table}

\begin{proposition}[The family of Sch\"onhage]\label{prop:shapes}
For each of the nine shapes of Table~\ref{tab:family} other than \((3,3)\),
the supremum over \(c>s+1\) of the saving \(S(c)\) of the shape lies in the
interval given in the table; in particular it is below \(0.00188\).  For
\((3,3)\) the saving exceeds \(\SstarLo\) at \(c=\cCond\).
\end{proposition}

\begin{proof}
As for Theorem~\ref{thm:global}, shape by shape: a lower bound is the value at
one rational ratio, and the upper bound is certified on cells covering
\((s+1,c_{\max}]\), with \(c_{\max}\) between \(200\) and \(600\), and beyond
\(c_{\max}\) by \(\Gam(c)<2/3\) and the monotonicity of \(\Apr\); the cells are
finest on the window around the maximum, which is \([19.6,20.6]\) for
\((2,4)\) and \([52.1,56.1]\) for \((5,5)\).
\end{proof}

In Lean: \lean{shapes_table} collects the nine upper bounds and the value for
\((3,3)\); the enclosures are \lean{sSup_shapeSaving_S24_mem} and its
analogues for the other shapes, from the per-shape bounds
\lean{shapeSaving_S24_lt} and \lean{shapeSaving_S24_gt}.  The
\LeanShapesCells{} cells, in \LeanShapesCellFiles{} generated files
\lean{Optimum/Shapes/Cells_S24_1.lean} and so on, are produced by
\texttt{search/certs.py} with the option \texttt{--encoding shapes} and
checked by Lean's kernel as in Theorem~\ref{thm:global}.

The shape \((3,3)\) is the best.  The exponent \(\Gam\) hardly varies over the
family, and \((3,3)\) also has the largest \(\varepsilon\).  For the same
reason mixtures do not help (\nf): when two shapes are used on different
levels, the optimizer drives the share of the second shape to zero, because no
trade-off between \(\Gam\) and \(\varepsilon\) is available.

\subsection{Hypothetical shapes}

This subsection is \nf; its numbers are floating point.  One can ask what an identity with \(s\) outer outputs, one term each, and a
shared term of inner length \(b\) would give, whether or not it exists.  The
saving grows with \(b\) and falls with \(s\).  For \(s=9\) it is \(0.00210\)
at \(b=4\), \(0.00254\) at \(b=5\) and \(0.00293\) at \(b=6\); for \(s=8\) it is
\(0.00188\) at \(b=3\) and \(0.00252\) at \(b=4\).  In the computed grid (\(5\le s\le12\), \(3\le b\le8\)) every cell with a
larger saving than \((s,b)=(9,4)\) is outside the class of Sch\"onhage's
construction: an identity of this kind with \(s=kn\) satisfies the pairing
bound \(b\le(k-1)(n-1)\), so \(s=9\) allows \(b\le4\), \(s=8=2\cdot4\) allows
\(b\le3\), \(s=6\) allows \(b\le2\), \(s=10\) allows \(b\le4\), \(s=12\) allows
\(b\le6\), and \(s=5,7\) have no factorization.  So within the class of
identities with one term per outer output, the identity of the paper is optimal
for this method.  Even the realistic cells beyond the bound
(\(s\in\{8,9\}\), \(b\in\{4,5\}\)) give savings of only \(0.0025\) to
\(0.0031\), that is, 3SUM in about \(n^{1.9985}\).

\subsection{Experiments on the arrangement of the wanted set}

This subsection is \nf.  The bound \eqref{eq:lemma11} is a worst case over the wanted sets.  A random
wanted set of the same size meets it up to lower-order factors: a leaf of
order \(d\) is shared by at most \(\binomial{L-m+d}{d}\) output strings, and the
expected number of visited leaves of order \(d\) is about
\(\min(|U|\alpha_d,\beta_d)\).  A gain needs a wanted set that is close to a
union of sub-tiles.  Simulation cannot test this: feasible sizes have
\(c=L/m\approx3\), where the \(\beta_d\) grow with \(d\) and every wanted set
saturates the tree, while the relevant regime \(c\approx21\) has at least
\(10^{21}\) leaves already for \(m=1\).  Zeros in the encodings do not help
either: an encoding at a leaf with \(\Pz\) on a set \(Z\) of levels is a signed
sum of at least \(3^{|Z|}\) entries, and for the dominant leaves
(\(|Z|\approx0.89m\)) the expected number of nonzero terms tends to infinity.

\subsection{Open problems}

\begin{enumerate}
\item \emph{New sharing structures.}  An identity in which outer outputs are
  fed by several terms, or with several inner outputs, is outside the leaf
  counting of Section~\ref{sec:recursion}.  Scoring such an identity needs a
  generalization of \(\alpha_d\) and \(\beta_d\).
\item \emph{Structured wanted sets.}  The reductions from 3SUM produce Exact
  Triangle instances whose weights are array entries at additively combined
  indices.  Whether the wanted sets are then far from random is not known.
\end{enumerate}

\begin{remark}[\nf]
The exponent of matrix multiplication plays no role.  The choice of the prime
costs \(n^{\omega}D'^{3/2}\), which is of lower order with Strassen's
\(\omega\le\log_2 7\) already; rectangular fast matrix multiplication would not
change any exponent of this paper.
\end{remark}


\subsection*{Acknowledgments}

The algorithms, the reductions and the framework of thin products are those
of Alman and Vassilevska Williams \cite{AV}; this paper changes their analysis
and their algorithm in a few places, and it would not exist without their work,
nor without the Lean formalization \cite{upstream} of their paper by Anthropic,
on which the development builds.  The proofs, the Lean development and the
text of this paper were written by the language model Claude Fable~5.1
(Anthropic), run as an agent in Claude Code, under the direction of the author,
who chose the goals, decided how the results are stated and reviewed the
outputs, and wrote no Lean by hand.  The repository
\url{https://github.com/qawbecrdtey/ImprovedExponents} contains the development
and this paper, and records the checks that were run
(Section~\ref{sec:formalization}).

\bibliographystyle{amsplain}
\bibliography{references}

\end{document}
